\documentclass[11pt]{article}

\usepackage[T1]{fontenc}
\usepackage[utf8]{inputenc}
\usepackage{lmodern}
\usepackage{microtype}
\usepackage[a4paper,margin=30mm,headheight=14pt]{geometry}
\usepackage{amsmath,amssymb,amsthm,mathtools,mathrsfs}
\usepackage{enumitem}
\usepackage{xcolor}
\usepackage{booktabs}
\usepackage{fancyhdr}
\usepackage{hyperref}

\hypersetup{
  colorlinks=true,
  linkcolor=blue!45!black,
  citecolor=blue!45!black,
  urlcolor=blue!45!black,
  pdftitle={Cubical Holonomy, Explicit SU(2) Flux, and a Finite Gaussian Cusp Limit},
  pdfauthor={Working draft}
}

\newtheorem{theorem}{Theorem}[section]
\newtheorem{proposition}[theorem]{Proposition}
\newtheorem{lemma}[theorem]{Lemma}
\newtheorem{corollary}[theorem]{Corollary}
\theoremstyle{definition}
\newtheorem{definition}[theorem]{Definition}
\newtheorem{remark}[theorem]{Remark}

\numberwithin{equation}{section}
\setlist{itemsep=2pt,topsep=4pt}
\setlength{\emergencystretch}{2em}
\linespread{1.025}
\pagestyle{fancy}
\fancyhf{}
\fancyhead[L]{Compact-link flux and finite Gaussian cusp limit}
\fancyhead[R]{Working draft}
\fancyfoot[C]{\thepage}

\title{\bfseries Cubical Holonomy, Explicit \(\mathrm{SU}(2)\) Flux,\\
and a Finite Gaussian Cusp Limit}
\author{Working draft}
\date{September 2026}

\begin{document}
\maketitle

\begin{abstract}
We give a self-contained finite construction on a family of flat four-tori
whose period matrices enter a controlled radial cusp.  Two integers are kept
independent: \(M_{\mathrm{cov}}\), which changes the covering lattice, and
\(N_{\mathrm{mesh}}\), which refines a transported cubical decomposition.
We identify the cubical vertices exactly with the finite torsion subgroup,
type the continuum-to-link holonomy map and its gauge covariance, and exhibit
a periodic diagonal \(\mathrm{SU}(2)\) flux.  Its logarithmic action is
\(2\pi^2M_{\mathrm{cov}}^2/(g_{\mathrm{YM}}^2D_T)\), while an explicitly
defined continuum-matched \((1,2)\)-sector Wilson functional has the same
mesh limit.  On the native fibre, a distinct smooth integral-flux
\(\mathrm{SU}(2)\) connection has the same numerical action but an exact
background-gauge Jacobi eigenvalue \(-4\pi/D_T\), so it is a classical saddle
rather than a vacuum; no equality with the displayed link holonomies is
asserted.  The scalar Laplacian on every cover has an exact character
spectrum; its deck-invariant sector is identified explicitly, and the cover
momentum lattices become dense when both cusp and cover parameters diverge.
A finite Coulomb-reduced cochain complex yields an exact Fourier spectrum and
a positive Gaussian oscillator gap at every fixed triple
\((T,M_{\mathrm{cov}},N_{\mathrm{mesh}})\).  Along
\(T_j=j^2\), \(M_{\mathrm{cov},j}=j\), and
\(N_{\mathrm{mesh},j}=j^3\), the cover tori converge pointed locally
isometrically to \(\mathbb R^4\), the mesh resolves the cover geometry, and
both the Gaussian gap and the flux action tend to zero.  The direct integral
of the mean-zero native scalar Laplacians over the radial parameter is a
specified self-adjoint operator with trivial kernel and zero in its essential
spectrum.  These statements
concern a finite quadratic oscillator and one explicitly normalized
compact-link sector.  On the native torus \(M_{\mathrm{cov}}=1\), we also
identify the exact reduced continuum Yang--Mills Hessian gap and the exact
finite gap in every mesh regime.  Finally, explicit gauge-invariant
probability measures show that the link kinematics and the displayed magnetic
configuration determine no Wilson-loop area law, while bounded self-adjoint
countermodels show that entanglement and a gapped Gaussian subspace determine
no full spectral gap.  We state and prove the full-form-domain coercive
comparison that would transfer an auxiliary gap to a physical Hamiltonian.
No such comparison is constructed here, so these results do not identify the
oscillator with a nonlinear, gauge-invariant compact-link Yang--Mills
Hamiltonian and do not constitute a physical mass-gap claim.
\end{abstract}

\medskip
\noindent\fcolorbox{black!55}{black!3}{%
\begin{minipage}{0.94\linewidth}
\textbf{Status and scope.}
This is a standalone mathematical draft of exact finite geometric,
compact-link, and Gaussian statements.  All period hypotheses and
normalizations used below are stated explicitly.  No continuum quantum field
theory is constructed; no transfer matrix, nonlinear Gauss-law Hamiltonian,
renormalized limit, or Osterwalder--Schrader reconstruction is asserted.
Consequently the Gaussian spectral gap proved here is not the spectral gap of
physical four-dimensional compact-link Yang--Mills theory.
\end{minipage}}

\section{Introduction and exact standing assumptions}

The purpose of this note is to separate, and then connect by explicit maps,
three finite constructions associated with a degenerating family of flat
four-tori.  The first is geometric: a cover is equipped with the cubical
decomposition transported from a coordinate torus.  The second is genuinely
compact-link: continuum connections produce group-valued edge holonomies,
and an explicit periodic \(\mathrm{SU}(2)\) flux has a computable plaquette
action.  The third is linear and Gaussian: a finite cochain complex on the
same mesh produces a positive Coulomb-restricted Hessian and a finite
harmonic oscillator.

The separation of parameters is essential.  The cover index
\(M_{\mathrm{cov}}\) multiplies two period vectors and changes the global
flat torus.  The mesh integer \(N_{\mathrm{mesh}}\) subdivides a fixed
coordinate torus and changes only its cubical resolution.  They are never
identified.  Only in the final diagonal sequence are both chosen as
functions of a third integer \(j\).

The period matrix and radial cusp specialization used below are taken from
\cite{AlpogeS6_2026}.  We restate every input needed here so that the proofs
depend only on the explicit regular-fibre matrix and the quantified radial
bounds below; neither the claimed compactification nor the claimed
identification of its total space with \(S^6\) is used.  Every cover below is
a fibrewise cover of the displayed regular flat torus.  No finite cover of a
compactified total space is asserted here.

Fix \(T_0>0\).  For every \(T\geq T_0\), let
\(\tau_T,\mu_T,\beta_T\in\mathbb C\), and put
\[
 L_T=\operatorname{Im}\tau_T,\qquad
 m_T=\operatorname{Im}\mu_T,\qquad
 q_T=-\operatorname{Im}\beta_T,\qquad
 D_T=L_Tq_T+6m_T^2.
\]
The standing period assumptions are the following exact quantified
requirements: \(D_T>0\) for every \(T\geq T_0\), and there is a constant
\(C_{\mathrm{per}}<\infty\), independent of \(T\), such that
\begin{equation}
 \label{exact-standing-period-assumptions}
 \begin{aligned}
 |L_T-T|&\leq C_{\mathrm{per}},&
 |q_T-T|&\leq C_{\mathrm{per}},&
 |m_T|&\leq C_{\mathrm{per}},\\
 |\operatorname{Re}\tau_T|&\leq C_{\mathrm{per}},&
 |\operatorname{Re}\mu_T|&\leq C_{\mathrm{per}},&
 |\operatorname{Re}\beta_T|&\leq C_{\mathrm{per}}.
 \end{aligned}
\end{equation}
Thus every \(O(1)\) below is uniform for \(T\geq T_0\), with a bound
depending only on the fixed period family.  We use the Euclidean metric on
\(\mathbb R^4\) and the induced quotient metrics.  The gauge coupling
\(g_{\mathrm{YM}}>0\) is fixed.

The ordered plaquette and real-trace structure used below is the classical
Wilson construction \cite{Wilson_1974}.  Transfer-matrix and reflection
positivity results such as \cite{Luscher_1977,MenottiPelissetto_1987} require
additional time/reflection structure, while Euclidean reconstruction
requires the axiomatic input of
\cite{OsterwalderSchrader_1973,OsterwalderSchrader_1975}.  These references
motivate the sharp boundary drawn in the last section; none is used to
replace a proof of the finite claims.

\subsection*{Results at a glance}

Starting from the exact period matrix in
\eqref{cusp-link-period-matrix}, we prove:

\begin{enumerate}
\item the scalar Laplacian on the cover has a complete character spectrum,
its deck-invariant Fourier sector is exact, and its momentum lattice has an
explicit covering-radius bound;
\item the transported cubulation covers the entire cover torus, and its
vertices are exactly, but only, the \(N_{\mathrm{mesh}}\)-torsion points;
\item continuum holonomy defines a typed gauge-equivariant map to link
variables, with the endpoint inverses fixed by the stated ODE convention;
\item an explicit periodic diagonal \(\mathrm{SU}(2)\) flux has constant
\((1,2)\)-plaquette and exact logarithmic and Wilson-sector actions;
\item on the native fibre, a separately constructed smooth integral-flux
connection has the same numerical action and a completely typed
background-gauge Jacobi operator with an explicit eigenvalue
\(-4\pi/D_T\); it is therefore a classical saddle, not a gapped vacuum;
\item the Coulomb-reduced finite Hessian has an exact Fourier spectrum, so
its Gaussian quantization has a unique ground state and explicit gap;
\item on the native torus, the reduced continuum Hessian and the finite
Hessian have exact first eigenvalues and multiplicities for every mesh size;
\item the direct integral of the native mean-zero scalar Laplacians over the
radial parameter has trivial kernel but zero in its essential spectrum;
\item one diagonal limit simultaneously decompactifies the cover, resolves
its mesh, and drives the finite Gaussian gap and flux action to zero; and
\item the exponential map is only a local chart, with an exact
\(N_{\mathrm{mesh}}^{-2}\) Hessian scaling, and supplies no global unitary
identification with compact-link gauge theory;
\item gauge-invariant pure-gauge and product-Haar states on the same link
configuration space give incompatible Wilson-loop expectations, so neither
the kinematics nor the displayed flux selects a confinement law; and
\item entanglement is not a spectral sector, entangled eigenvalues can
accumulate at zero, and an embedded gapped Gaussian range can coexist with a
gapless orthogonal channel; a physical gap transfer requires a comparison on
the entire physical excitation form domain.
\end{enumerate}

\section{Cubical holonomy and finite gauge systems}
\label{cusp-link-section}

This section records three constructions on the regular cusp fibres while
keeping two independent integers separate throughout:
\(M_{\mathrm{cov}}\) is the index used to enlarge two periods of a fibre,
whereas \(N_{\mathrm{mesh}}\) is the number of cubical subdivisions in each
coordinate direction.  In particular, no mesh refinement is identified with
a change of covering space.

Let
\begin{equation}
 P_T=
 \begin{pmatrix}
  6\operatorname{Re}\mu_T&\operatorname{Re}\tau_T&1&0\\
  6\operatorname{Im}\mu_T&\operatorname{Im}\tau_T&0&0\\
  \operatorname{Re}\beta_T&\operatorname{Re}\mu_T&0&1\\
  \operatorname{Im}\beta_T&\operatorname{Im}\mu_T&0&0
 \end{pmatrix},
 \qquad
 F_T=\mathbb R^4/P_T\mathbb Z^4.
 \label{cusp-link-period-matrix}
\end{equation}
Write
\begin{equation}
 L_T=\operatorname{Im}\tau_T,
 \qquad m_T=\operatorname{Im}\mu_T,
 \qquad q_T=-\operatorname{Im}\beta_T,
 \qquad D_T=L_Tq_T+6m_T^2>0.
 \label{cusp-link-radial-data}
\end{equation}
Expanding the determinant along the last two columns gives
\[
 \det P_T=L_T\operatorname{Im}\beta_T-6m_T^2=-D_T.
\]
Along the positive radial cusp the source data give
\begin{equation}
 L_T=T+O(1),\qquad q_T=T+O(1),\qquad m_T=O(1),
 \label{cusp-link-cusp-asymptotics}
\end{equation}
and all three real parts occurring in \eqref{cusp-link-period-matrix} are
bounded.  Equivalently, after ordering physical coordinates as
\((x_1,x_3;x_2,x_4)\), define the upper block
\[
 A_T=
 \begin{pmatrix}
  6\operatorname{Re}\mu_T&\operatorname{Re}\tau_T\\
  \operatorname{Re}\beta_T&\operatorname{Re}\mu_T
 \end{pmatrix}.
\]
The first two columns then have block form
\begin{equation}
 \binom{A_T}{B_T},\qquad
 B_T=
 \begin{pmatrix}6m_T&L_T\\-q_T&m_T\end{pmatrix}
 =T\begin{pmatrix}0&1\\-1&0\end{pmatrix}+C_T,
 \qquad \sup_T\lVert C_T\rVert_{\mathrm{op}}<\infty.
 \label{cusp-link-block-form}
\end{equation}

For \(M_{\mathrm{cov}}\geq1\), put
\begin{equation}
 S_{M_{\mathrm{cov}}}
 =\operatorname{diag}(1,1,M_{\mathrm{cov}},M_{\mathrm{cov}}),
 \qquad
 P_{T,M_{\mathrm{cov}}}=P_TS_{M_{\mathrm{cov}}},
 \label{cusp-link-cover-matrix}
\end{equation}
and define the degree-\(M_{\mathrm{cov}}^2\) cover
\begin{equation}
 \widetilde F_{T,M_{\mathrm{cov}}}
 =\mathbb R^4/P_{T,M_{\mathrm{cov}}}\mathbb Z^4
 \longrightarrow F_T.
 \label{cusp-link-cover-torus}
\end{equation}
Indeed,
\(S_{M_{\mathrm{cov}}}\mathbb Z^4\subset\mathbb Z^4\) and
\[
 \mathbb Z^4/S_{M_{\mathrm{cov}}}\mathbb Z^4
 \cong(\mathbb Z/M_{\mathrm{cov}}\mathbb Z)^2,
\]
so the identity on \(\mathbb R^4\) induces the asserted cover and its deck
group has order \(M_{\mathrm{cov}}^2\).
Its Euclidean quotient metric, written on the coordinate torus
\(\mathbb R^4/\mathbb Z^4\), has Gram matrix and covolume
\begin{equation}
 G_{T,M_{\mathrm{cov}}}
 =P_{T,M_{\mathrm{cov}}}^{\mathsf T}P_{T,M_{\mathrm{cov}}},
 \qquad
 \mathcal D_{T,M_{\mathrm{cov}}}
 =\sqrt{\det G_{T,M_{\mathrm{cov}}}}
 =M_{\mathrm{cov}}^2D_T.
 \label{cusp-link-cover-metric}
\end{equation}

\subsection{Scalar characters and exact cover momentum sectors}

Equip \(\widetilde F_{T,M_{\mathrm{cov}}}\) with normalized Haar measure
\begin{equation}
 \mathrm d\nu_{T,M_{\mathrm{cov}}}
 =\mathcal D_{T,M_{\mathrm{cov}}}^{-1}\,
   \mathrm d\operatorname{vol}_{\widetilde F_{T,M_{\mathrm{cov}}}}.
 \label{cusp-link-normalized-cover-measure}
\end{equation}
Let \(\widetilde\Delta_{T,M_{\mathrm{cov}}}\) be the nonnegative
self-adjoint scalar Laplacian obtained from the closed form
\begin{equation}
 \widetilde{\mathfrak q}_{T,M_{\mathrm{cov}}}[f]
 =\int_{\widetilde F_{T,M_{\mathrm{cov}}}}|\mathrm df|^2\,
       \mathrm d\nu_{T,M_{\mathrm{cov}}},
 \qquad
 \operatorname{Dom}(\widetilde{\mathfrak q}_{T,M_{\mathrm{cov}}})=H^1,
 \label{cusp-link-cover-dirichlet-form}
\end{equation}
whose operator domain is \(H^2(\widetilde F_{T,M_{\mathrm{cov}}})\).

\begin{theorem}[Cover spectrum, invariant sector, and momentum density]
\label{cusp-link-cover-spectrum}
The functions
\begin{equation}
 \widetilde\chi_{T,M_{\mathrm{cov}},k}([x])
 =\exp\!\left(
   2\pi i\left\langle
    P_T^{-\mathsf T}S_{M_{\mathrm{cov}}}^{-1}k,x
   \right\rangle\right),
 \qquad k\in\mathbb Z^4,
 \label{cusp-link-cover-characters}
\end{equation}
form a complete orthonormal basis, and
\begin{equation}
 \widetilde\Delta_{T,M_{\mathrm{cov}}}
 \widetilde\chi_{T,M_{\mathrm{cov}},k}
 =4\pi^2
  \left|P_T^{-\mathsf T}S_{M_{\mathrm{cov}}}^{-1}k\right|^2
  \widetilde\chi_{T,M_{\mathrm{cov}},k}.
 \label{cusp-link-cover-scalar-eigenvalue}
\end{equation}
If \(c_{T,M_{\mathrm{cov}}}:\widetilde F_{T,M_{\mathrm{cov}}}\to F_T\)
is the covering map and
\(\chi_{T,\ell}([x])=\exp(2\pi i\langle P_T^{-\mathsf T}\ell,x\rangle)\),
then
\begin{equation}
 c_{T,M_{\mathrm{cov}}}^*\chi_{T,\ell}
 =\widetilde\chi_{T,M_{\mathrm{cov}},S_{M_{\mathrm{cov}}}\ell}.
 \label{cusp-link-cover-character-pullback}
\end{equation}
Consequently, the deck-invariant sector is exactly
\begin{equation}
 k\in S_{M_{\mathrm{cov}}}\mathbb Z^4
 \quad\Longleftrightarrow\quad
 k_3\equiv k_4\equiv0\pmod {M_{\mathrm{cov}}}.
 \label{cusp-link-cover-deck-sector}
\end{equation}
There is a constant \(C_{\mathrm{mom}}\), independent of all sufficiently
large \(T\) and all \(M_{\mathrm{cov}}\ge1\), such that
\begin{equation}
 \sup_{p\in\mathbb R^4}\operatorname{dist}\!\left(
  p,P_T^{-\mathsf T}S_{M_{\mathrm{cov}}}^{-1}\mathbb Z^4
 \right)
 \le C_{\mathrm{mom}}\left(\frac1T+\frac1{M_{\mathrm{cov}}}\right).
 \label{cusp-link-cover-momentum-radius}
\end{equation}
Thus, whenever \(T_j,M_{\mathrm{cov},j}\to\infty\), these momentum lattices
become locally dense in \(\mathbb R^4\), and every \(E\ge0\) is a limit of
eigenvalues in \eqref{cusp-link-cover-scalar-eigenvalue}.
\end{theorem}

\begin{proof}
The dual lattice of
\(P_TS_{M_{\mathrm{cov}}}\mathbb Z^4\) is
\begin{equation}
 (P_TS_{M_{\mathrm{cov}}})^{-\mathsf T}\mathbb Z^4
 =P_T^{-\mathsf T}S_{M_{\mathrm{cov}}}^{-1}\mathbb Z^4.
 \label{cusp-link-cover-dual-lattice}
\end{equation}
A character of \(\mathbb R^4\) descends to the quotient precisely when its
frequency lies in this dual lattice.  Fourier theory for the compact abelian
group, normalized by \eqref{cusp-link-normalized-cover-measure}, gives the
complete orthonormal basis \eqref{cusp-link-cover-characters}; applying the
positive Euclidean Laplacian proves
\eqref{cusp-link-cover-scalar-eigenvalue}.  Substituting
\(k=S_{M_{\mathrm{cov}}}\ell\) proves
\eqref{cusp-link-cover-character-pullback}.  A cover character descends
exactly when it is deck invariant, so
\eqref{cusp-link-cover-character-pullback} and Fourier completeness prove
\eqref{cusp-link-cover-deck-sector}.

For the quantitative assertion, abbreviate
\[
 \alpha_T=\operatorname{Re}\mu_T,\qquad
 \eta_T=\operatorname{Re}\tau_T,\qquad
 c_T=\operatorname{Re}\beta_T,\qquad M=M_{\mathrm{cov}}.
\]
Solving \(P_T^{\mathsf T}x=e_j\) gives the four columns
\begin{align*}
 b_1&=\left(0,\frac{m_T}{D_T},0,-\frac{L_T}{D_T}\right),\\
 b_2&=\left(0,\frac{q_T}{D_T},0,\frac{6m_T}{D_T}\right),\\
 b_3&=\frac1M\left(
 1,\frac{-6m_T\alpha_T-q_T\eta_T}{D_T},0,
 \frac{6(L_T\alpha_T-m_T\eta_T)}{D_T}\right),\\
 b_4&=\frac1M\left(
 0,\frac{-m_Tc_T-q_T\alpha_T}{D_T},1,
 \frac{L_Tc_T-6m_T\alpha_T}{D_T}\right)
\end{align*}
of \(P_T^{-\mathsf T}S_M^{-1}\).  The standing bounds imply
\(D_T\asymp T^2\) and hence, for a uniform \(C_0\),
\[
 |b_1|+|b_2|\le \frac{C_0}{T},
 \qquad |b_3|+|b_4|\le \frac{C_0}{M}.
\]
Write any \(p\in\mathbb R^4\) uniquely as \(p=\sum_jt_jb_j\), and choose
\(n_j\in\mathbb Z\) with \(|t_j-n_j|\le\tfrac12\).  Then
\[
 \left|p-\sum_jn_jb_j\right|
 \le\frac12\sum_j|b_j|
 \le\frac{C_0}{2}\left(\frac1T+\frac1M\right),
\]
which proves \eqref{cusp-link-cover-momentum-radius}.  If \(E>0\), choose
\(p\) with \(4\pi^2|p|^2=E\) and approximate it by the cover momentum
lattice; continuity gives eigenvalues tending to \(E\).  For \(E=0\), the
index \(k=(1,0,0,0)^{\mathsf T}\) has momentum norm \(O(T^{-1})\).
\end{proof}

\begin{remark}[Fibrewise scope]
Theorem~\ref{cusp-link-cover-spectrum} is entirely fibrewise.  It does not
assert that these covers assemble compatibly with monodromy over a punctured
base or extend across a singular fibre.
\end{remark}

\begin{proposition}[Native and expanding-cover pointed limits]
\label{cusp-link-pointed-limits}
Let
\[
 \mathbb T^2_{\mathrm{re}}
 = (\mathbb Re_1\oplus\mathbb Re_3)/
   (\mathbb Ze_1\oplus\mathbb Ze_3).
\]
Then the pointed native tori satisfy
\begin{equation}
 (F_T,0)\longrightarrow
 (\mathbb T^2_{\mathrm{re}}\times
   (\mathbb Re_2\oplus\mathbb Re_4),(0,0))
 \cong(\mathbb T^2\times\mathbb R^2,(0,0)),
 \label{cusp-link-native-pointed-limit}
\end{equation}
in pointed smooth topology: each fixed-radius ball is eventually exactly
isometric to the corresponding ball in the limit.  There are constants
\(C,C'>0\) such that, for all sufficiently large \(T\) and every
\(M\ge1\),
\begin{equation}
 \min\{T-C,M\}
 \le\operatorname{sys}(\widetilde F_{T,M})
 \le\min\{T+C',M\}.
 \label{cusp-link-cover-systole-two-sided}
\end{equation}
Consequently, for every integer-valued \(M(T)\to\infty\),
\begin{equation}
 (\widetilde F_{T,M(T)},0)\longrightarrow(\mathbb R^4,0)
 \label{cusp-link-cover-pointed-limit}
\end{equation}
with the same exact fixed-ball property.  In particular, the native and
expanding-cover unscaled pointed limits are different.
\end{proposition}

\begin{proof}
The period lattice of \(F_T\) is generated by the first two columns of
\(P_T\), denoted \(v_{\gamma,T},v_{u,T}\), together with \(e_1,e_3\).
First quotient by \(\mathbb Ze_1\oplus\mathbb Ze_3\).  The projection of
\(av_{\gamma,T}+bv_{u,T}\) to the noncompact
\((e_2,e_4)\)-plane is
\[
 B_T\binom ab
 =\left(T\begin{pmatrix}0&1\\-1&0\end{pmatrix}+C_T\right)\binom ab.
\]
With \(C=\sup_T\|C_T\|_{\mathrm{op}}\), its norm is at least
\((T-C)\sqrt{a^2+b^2}\).  Hence the remaining nontrivial quotient
translations eventually displace every point by more than \(4R\).  The
quotient map is then injective and distance-preserving on the radius-\(R\)
ball, proving \eqref{cusp-link-native-pointed-limit}.

Every nonzero vector in the cover lattice has the form
\[
 av_{\gamma,T}+bv_{u,T}+M(pe_1+qe_3),
 \qquad a,b,p,q\in\mathbb Z.
\]
If \((a,b)\ne(0,0)\), its noncompact-plane projection has norm at least
\(T-C\); if \((a,b)=(0,0)\), its norm is at least \(M\).  This gives the
lower bound in \eqref{cusp-link-cover-systole-two-sided}.  The vector
\(Me_1\) has length \(M\), while either long generator has length at most
\(T+C'\) for a uniform \(C'\), proving the upper bound.  If both \(T\) and
\(M(T)\) diverge, the systole diverges.  Once it exceeds \(4R\), the universal
covering map identifies the radius-\(R\) balls isometrically, proving
\eqref{cusp-link-cover-pointed-limit}.
\end{proof}

\subsection{The transported cubical subdivision and continuum holonomy}

Fix an integer \(N_{\mathrm{mesh}}\geq2\), and let
\(K_{N_{\mathrm{mesh}}}\) be the standard cubical subdivision of
\(\mathbb R^4/\mathbb Z^4\) obtained by dividing every coordinate circle
into \(N_{\mathrm{mesh}}\) equal intervals.  Its vertex set and positively
oriented coordinate edges are
\begin{equation}
 V_{N_{\mathrm{mesh}}}
 =(\mathbb Z/N_{\mathrm{mesh}}\mathbb Z)^4,
 \qquad
 E_{N_{\mathrm{mesh}}}
 =V_{N_{\mathrm{mesh}}}\times\{1,2,3,4\}.
 \label{cusp-link-cubical-sets}
\end{equation}
Here and below, \(\mathbf e_1,\ldots,\mathbf e_4\) denotes the standard
basis of \(\mathbb R^4\), reduced modulo \(N_{\mathrm{mesh}}\) when used to
translate an element of \(V_{N_{\mathrm{mesh}}}\).

\begin{proposition}[Exact cubical realization and the torsion vertices]
\label{cusp-link-cubical-realization}
The linear map \(P_{T,M_{\mathrm{cov}}}\) induces a diffeomorphism
\begin{equation}
 \Phi_{T,M_{\mathrm{cov}}}:\mathbb R^4/\mathbb Z^4
 \longrightarrow\widetilde F_{T,M_{\mathrm{cov}}},
 \qquad [y]\longmapsto[P_{T,M_{\mathrm{cov}}}y].
 \label{cusp-link-coordinate-diffeomorphism}
\end{equation}
Transporting \(K_{N_{\mathrm{mesh}}}\) by this diffeomorphism gives an
honest cubical cell decomposition of the whole torus.  On vertices its
restriction is the bijection
\begin{equation}
 \begin{aligned}
 \iota_{T,M_{\mathrm{cov}},N_{\mathrm{mesh}}}:\quad
 &(\mathbb Z/N_{\mathrm{mesh}}\mathbb Z)^4
 \xrightarrow{\ \cong\ }
 \widetilde F_{T,M_{\mathrm{cov}}}[N_{\mathrm{mesh}}],\\
 &[n]\longmapsto
 \left[\frac1{N_{\mathrm{mesh}}}P_{T,M_{\mathrm{cov}}}n\right].
 \end{aligned}
 \label{cusp-link-torsion-bijection}
\end{equation}
Here the codomain is the finite \(N_{\mathrm{mesh}}\)-torsion subgroup,
not the entire torus.  The coordinate edge \((n,j)\) is the path
\begin{equation}
 \gamma_{n,j}(t)=
 \left[
  \frac1{N_{\mathrm{mesh}}}
  P_{T,M_{\mathrm{cov}}}(n+t\mathbf e_j)
 \right],
 \qquad 0\leq t\leq1.
 \label{cusp-link-coordinate-edge}
\end{equation}
More generally, for an ordered subset
\(I=\{i_1,\ldots,i_p\}\subset\{1,2,3,4\}\), the corresponding closed
\(p\)-cell is parametrized by
\begin{equation}
 c_{n,I}(t_1,\ldots,t_p)=
 \left[
  \frac1{N_{\mathrm{mesh}}}P_{T,M_{\mathrm{cov}}}
  \left(n+\sum_{a=1}^pt_a\mathbf e_{i_a}\right)
 \right],
 \qquad (t_1,\ldots,t_p)\in[0,1]^p.
 \label{cusp-link-general-cell}
\end{equation}
\end{proposition}

\begin{proof}
If \(y-y'\in\mathbb Z^4\), then
\(P_{T,M_{\mathrm{cov}}}(y-y')\) belongs to the quotient lattice, so
\eqref{cusp-link-coordinate-diffeomorphism} is well defined.  The matrix
\(P_{T,M_{\mathrm{cov}}}\) is invertible, and its inverse induces the
inverse map on the two quotients.  Hence it is a diffeomorphism and carries
every open cell of the standard subdivision diffeomorphically onto an open
cell of the target.

Every point displayed in \eqref{cusp-link-torsion-bijection} is killed by
\(N_{\mathrm{mesh}}\).  If it is zero, then
\[
 \frac{P_{T,M_{\mathrm{cov}}}n}{N_{\mathrm{mesh}}}
 \in P_{T,M_{\mathrm{cov}}}\mathbb Z^4.
\]
Invertibility gives
\(n\in N_{\mathrm{mesh}}\mathbb Z^4\), proving injectivity.  Conversely,
if \([x]\) is killed by \(N_{\mathrm{mesh}}\), choose a lift
\(x\in\mathbb R^4\).  There is \(n\in\mathbb Z^4\) such that
\(N_{\mathrm{mesh}}x=P_{T,M_{\mathrm{cov}}}n\), and therefore
\([x]=[P_{T,M_{\mathrm{cov}}}n/N_{\mathrm{mesh}}]\).  This proves
surjectivity.  Formula \eqref{cusp-link-coordinate-edge} is simply the
transport of the corresponding coordinate edge; at a wrap-around endpoint
the extra vector \(P_{T,M_{\mathrm{cov}}}\mathbf e_j\) is a period.
\end{proof}

Let \(G\) be a compact matrix Lie group and use the trivial principal bundle over
\(\widetilde F_{T,M_{\mathrm{cov}}}\).  A connection is represented by
\(A\in\Omega^1(\widetilde F_{T,M_{\mathrm{cov}}};\mathfrak g)\).  For an
oriented path \(\gamma\), define \(U_A(\gamma)\) by
\begin{equation}
 \frac{dU}{dt}=-A_{\gamma(t)}(\dot\gamma(t))U(t),
 \qquad U(0)=1,
 \qquad U_A(\gamma)=U(1).
 \label{cusp-link-parallel-transport-ode}
\end{equation}

\begin{proposition}[Typed continuum-to-link holonomy map]
\label{cusp-link-holonomy-map}
The formula
\begin{equation}
 \operatorname{Hol}_{T,M_{\mathrm{cov}},N_{\mathrm{mesh}}}(A)_{n,j}
 =U_A(\gamma_{n,j})
 \label{cusp-link-holonomy-definition}
\end{equation}
defines a map
\begin{equation}
 \operatorname{Hol}_{T,M_{\mathrm{cov}},N_{\mathrm{mesh}}}:
 \Omega^1(\widetilde F_{T,M_{\mathrm{cov}}};\mathfrak g)
 \longrightarrow G^{E_{N_{\mathrm{mesh}}}}.
 \label{cusp-link-holonomy-typed}
\end{equation}
For the convention
\begin{equation}
 A^g=g^{-1}Ag+g^{-1}dg,
 \qquad
 g\in C^\infty(\widetilde F_{T,M_{\mathrm{cov}}},G),
 \label{cusp-link-gauge-convention}
\end{equation}
it is equivariant for evaluation of \(g\) at the vertices:
\begin{equation}
 \operatorname{Hol}(A^g)_{n,j}
 =g(\gamma_{n,j}(1))^{-1}
  \operatorname{Hol}(A)_{n,j}
  g(\gamma_{n,j}(0)).
 \label{cusp-link-holonomy-equivariance}
\end{equation}
\end{proposition}

\begin{proof}
The smooth matrix ODE \eqref{cusp-link-parallel-transport-ode} has a unique
solution on every edge, which proves that \eqref{cusp-link-holonomy-typed}
is defined.  Let \(U\) solve that ODE for \(A\), set
\(g_t=g(\gamma_{n,j}(t))\), and put
\(V(t)=g_t^{-1}U(t)g_0\).  Direct differentiation gives
\[
 V'(t)=-\bigl(g_t^{-1}A(\dot\gamma_{n,j})g_t
                 +g_t^{-1}g_t'\bigr)V(t),
 \qquad V(0)=1.
\]
This is exactly the parallel-transport equation for \(A^g\).  Uniqueness
and evaluation at \(t=1\) prove \eqref{cusp-link-holonomy-equivariance}.
Notice that the endpoint inverse in this formula is forced by the joint
conventions \eqref{cusp-link-parallel-transport-ode} and
\eqref{cusp-link-gauge-convention}.
\end{proof}

\subsection{An explicit compact-link \texorpdfstring{\(\mathrm{SU}(2)\)}{SU(2)} flux}

The ordered plaquette convention used below is most naturally written for
the inverse-transport link variable
\begin{equation}
 \mathcal U(A)_{n,j}
 =\operatorname{Hol}_{T,M_{\mathrm{cov}},N_{\mathrm{mesh}}}(A)_{n,j}^{-1}.
 \label{cusp-link-inverse-transport-variable}
\end{equation}
Equation \eqref{cusp-link-holonomy-equivariance} then becomes
\begin{equation}
 \mathcal U(A^g)_{n,j}
 =g(\gamma_{n,j}(0))^{-1}\mathcal U(A)_{n,j}g(\gamma_{n,j}(1)).
 \label{cusp-link-inverse-transport-gauge-law}
\end{equation}
Thus the product in \eqref{cusp-link-plaquette-convention} transforms by
conjugation at its base vertex.  The symbols \(U_j(n)\) in the remainder of
this subsection use this inverse-transport convention.

Use
\begin{equation}
 \mathfrak{su}(2)=\{X\in M_2(\mathbb C):X^*=-X,\ \operatorname{tr}X=0\},
 \qquad
 \langle X,Y\rangle_{\mathfrak{su}(2)}
 =-\frac12\operatorname{tr}(XY),
 \label{cusp-link-su2-inner-product}
\end{equation}
on real elements, and fix
\begin{equation}
 H_{\mathrm{link}}=\begin{pmatrix}i&0\\0&-i\end{pmatrix},
 \qquad
 \lVert H_{\mathrm{link}}\rVert^2=1,
 \qquad \exp(2\pi H_{\mathrm{link}})=1.
 \label{cusp-link-flux-generator}
\end{equation}
Write \(N=N_{\mathrm{mesh}}\) temporarily inside the next displayed link
formulas, and represent \(n_1,n_2\) by integers in
\(\{0,\ldots,N-1\}\).

The link field constructed next is an abstract compact-link configuration.
It is not asserted to lie in the image of the smooth-connection holonomy map
\eqref{cusp-link-holonomy-typed}.

\begin{theorem}[A periodic compact-link flux and its exact action]
\label{cusp-link-flux-theorem}
On the oriented edges of \(K_N\), define
\begin{align}
 U_1(n)&=\exp\left(-\frac{2\pi n_2}{N^2}H_{\mathrm{link}}\right),
 \label{cusp-link-flux-u1}\\
 U_2(n)&=
 \begin{cases}
  1,&0\leq n_2\leq N-2,\\[1mm]
  \exp(2\pi n_1H_{\mathrm{link}}/N),&n_2=N-1,
 \end{cases}
 \label{cusp-link-flux-u2}\\
 U_3(n)&=U_4(n)=1.
 \label{cusp-link-flux-u34}
\end{align}
For
\begin{equation}
 U_{ij}(n)=U_i(n)U_j(n+\mathbf e_i)
            U_i(n+\mathbf e_j)^{-1}U_j(n)^{-1},
 \label{cusp-link-plaquette-convention}
\end{equation}
one has
\begin{equation}
 \begin{aligned}
 U_{12}(n)&=\exp(2\pi H_{\mathrm{link}}/N^2)
 &&\text{for every }n,\\
 U_{ij}(n)&=1
 &&\text{for }1\leq i<j\leq4,\quad(i,j)\neq(1,2).
 \end{aligned}
 \label{cusp-link-plaquette-values}
\end{equation}

Let \(\operatorname{Log}\) be the principal logarithm on the displayed
plaquette.  It is unambiguous because \(0<2\pi/N^2\leq\pi/2\).  Define
the logarithmic curvature two-form and its metric-weighted action by
\begin{align}
 \mathcal F_N^{\log}(n)
 &=N^2\operatorname{Log}(U_{12}(n))\,dy^1\wedge dy^2,
 \label{cusp-link-log-curvature}\\
 S^{\log}_{T,M_{\mathrm{cov}},N}(U)
 &=\frac{\mathcal D_{T,M_{\mathrm{cov}}}}{2g_{\mathrm{YM}}^2N^4}
   \sum_{n\in V_N}
   \lVert\mathcal F_N^{\log}(n)\rVert^2_{
       G_{T,M_{\mathrm{cov}}}^{-1}\otimes\mathfrak{su}(2)}.
 \label{cusp-link-log-action}
\end{align}
Then
\begin{equation}
 \mathcal F_N^{\log}(n)=2\pi H_{\mathrm{link}}\,dy^1\wedge dy^2,
 \qquad
 S^{\log}_{T,M_{\mathrm{cov}},N}(U)
 =\frac{2\pi^2M_{\mathrm{cov}}^2}{g_{\mathrm{YM}}^2D_T}.
 \label{cusp-link-log-action-value}
\end{equation}
In particular, the native-fibre value \(M_{\mathrm{cov}}=1\) is
\(2\pi^2/(g_{\mathrm{YM}}^2D_T)\), independently of the mesh.

There is no canonical scalar Wilson action for a general nonorthogonal
metric obtained merely by assigning independent character weights to all
coordinate plaquettes.  For this configuration, however, the exact
continuum-matched \((1,2)\)-sector Wilson functional is unambiguous if it is
defined explicitly by the following formula.  Its ordered plaquette
holonomy and real trace are the Wilson lattice-action construction
\cite{Wilson_1974}; the metric coefficient and normalization here are the
present definition, and their consequences are proved below rather than
imported from that source:
\begin{equation}
 S^{W,12}_{T,M_{\mathrm{cov}},N}(U)
 =\frac{\mathcal D_{T,M_{\mathrm{cov}}}\kappa^{12}_{T,M_{\mathrm{cov}}}}
        {2g_{\mathrm{YM}}^2}
   \sum_{n\in V_N}
   \left(2-\operatorname{Re}\operatorname{tr}U_{12}(n)\right),
 \label{cusp-link-wilson-sector-definition}
\end{equation}
where
\begin{equation}
 \kappa^{12}_{T,M_{\mathrm{cov}}}
 =(G_{T,M_{\mathrm{cov}}}^{-1})_{11}
  (G_{T,M_{\mathrm{cov}}}^{-1})_{22}
 -(G_{T,M_{\mathrm{cov}}}^{-1})_{12}^2
 =D_T^{-2}.
 \label{cusp-link-two-form-minor}
\end{equation}
Its exact value and continuum-mesh limit are
\begin{align}
 S^{W,12}_{T,M_{\mathrm{cov}},N}(U)
 &=\frac{M_{\mathrm{cov}}^2N^4}{2g_{\mathrm{YM}}^2D_T}
   \left(2-2\cos\frac{2\pi}{N^2}\right),
 \label{cusp-link-wilson-sector-value}\\
 &=\frac{M_{\mathrm{cov}}^2}{g_{\mathrm{YM}}^2D_T}
   \left(
    2\pi^2-\frac{2\pi^4}{3N^4}
    +\frac{4\pi^6}{45N^8}+O(N^{-12})
   \right),
 \label{cusp-link-wilson-sector-expansion}\\
 \lim_{N\to\infty}S^{W,12}_{T,M_{\mathrm{cov}},N}(U)
 &=\frac{2\pi^2M_{\mathrm{cov}}^2}{g_{\mathrm{YM}}^2D_T}.
 \label{cusp-link-wilson-sector-limit}
\end{align}
\end{theorem}

\begin{proof}
All links lie in the same diagonal circle subgroup, so they commute.  If
\(0\leq n_2\leq N-2\), the two \(U_2\)-terms in
\eqref{cusp-link-plaquette-convention} are trivial and
\[
 U_{12}(n)=
 \exp(-2\pi n_2H_{\mathrm{link}}/N^2)
 \exp(2\pi(n_2+1)H_{\mathrm{link}}/N^2)
 =\exp(2\pi H_{\mathrm{link}}/N^2).
\]
If \(n_2=N-1\), then \(U_1(n+\mathbf e_2)=1\).  Moreover,
\[
 U_2(n+\mathbf e_1)U_2(n)^{-1}=\exp(2\pi H_{\mathrm{link}}/N).
\]
For \(n_1=N-1\), the exponent computed with representatives differs from
\(2\pi H_{\mathrm{link}}/N\) by \(-2\pi H_{\mathrm{link}}\), whose exponential is one; hence the formula
also holds across that wrap boundary.  It follows that
\[
 U_{12}(n)=
 \exp\left(-\frac{2\pi(N-1)}{N^2}H_{\mathrm{link}}\right)
 \exp\left(\frac{2\pi}{N}H_{\mathrm{link}}\right)
 =\exp(2\pi H_{\mathrm{link}}/N^2).
\]
The functions \(U_1,U_2\) are independent of \(n_3,n_4\), and
\(U_3=U_4=1\), proving every remaining assertion in
\eqref{cusp-link-plaquette-values}.  Reversing the orientation gives
\(U_{21}(n)=U_{12}(n)^{-1}\); it is not an additional trivial plaquette.

The lower-right \(2\times2\) principal block of the cover Gram matrix is
\(M_{\mathrm{cov}}^2I_2\).  Jacobi's
complementary-minor identity therefore gives
\[
 \kappa^{12}_{T,M_{\mathrm{cov}}}
 =\frac{\det(M_{\mathrm{cov}}^2I_2)}
        {\det G_{T,M_{\mathrm{cov}}}}
 =\frac{M_{\mathrm{cov}}^4}{M_{\mathrm{cov}}^4D_T^2}
 =D_T^{-2},
\]
which proves \eqref{cusp-link-two-form-minor}.  By
\eqref{cusp-link-flux-generator} and the principal-log choice,
\(N^2\operatorname{Log}U_{12}=2\pi H_{\mathrm{link}}\).  Hence every one of the \(N^4\)
summands in \eqref{cusp-link-log-action} equals
\(4\pi^2D_T^{-2}\).  Substitution of
\(\mathcal D_{T,M_{\mathrm{cov}}}=M_{\mathrm{cov}}^2D_T\) proves
\eqref{cusp-link-log-action-value}.

Finally, \(H_{\mathrm{link}}\) has eigenvalues \(\pm i\), so
\[
 \operatorname{Re}\operatorname{tr}\exp(2\pi H_{\mathrm{link}}/N^2)
 =2\cos(2\pi/N^2).
\]
There are \(N^4\) equal plaquettes.  Substitution proves
\eqref{cusp-link-wilson-sector-value}.  Applying
\[
 1-\cos x=\frac{x^2}{2}-\frac{x^4}{24}
           +\frac{x^6}{720}+O(x^8)
\]
with \(x=2\pi/N^2\) proves
\eqref{cusp-link-wilson-sector-expansion} and
\eqref{cusp-link-wilson-sector-limit}.
\end{proof}

\subsection{A distinct smooth native magnetic saddle}
\label{cusp-link-smooth-magnetic-saddle-section}

The compact-link field above was defined directly on edges.  We now construct
a smooth integral-flux connection on the native torus \(F_T\).  Its action
equals the native value in \eqref{cusp-link-log-action-value}, but equality of
two real numbers is not an identification of connections: no claim is made
that the links \eqref{cusp-link-flux-u1}--\eqref{cusp-link-flux-u34} are the
holonomies of the connection below.

Use the physical coordinate order \(x=(x_1,y_1,x_2,y_2)^{\mathsf T}\), put
\(u=P_T^{-1}x\), and define the integral coframe
\begin{equation}
 \theta_T^j=du_j=(P_T^{-\mathsf T}\mathbf e_j)^{\mathsf T}dx.
 \label{cusp-link-magnetic-integral-coframe}
\end{equation}
For the marked cycle \(s\mapsto[x+sP_T\mathbf e_k]\), one has
\(\int\theta_T^j=\delta_{jk}\).  Direct solution of
\(P_T^{\mathsf T}\xi=\mathbf e_1,\mathbf e_2\), without replacing the skew
lattice by a rectangle, gives
\begin{equation}
 \theta_T^1=\frac{m_T}{D_T}\,dy_1-\frac{L_T}{D_T}\,dy_2,
 \qquad
 \theta_T^2=\frac{q_T}{D_T}\,dy_1+\frac{6m_T}{D_T}\,dy_2.
 \label{cusp-link-magnetic-first-coframe}
\end{equation}
Consequently,
\begin{equation}
 \theta_T^1\wedge\theta_T^2
 =\frac1{D_T}\,dy_1\wedge dy_2,
 \qquad
 |\theta_T^1\wedge\theta_T^2|^2=D_T^{-2}.
 \label{cusp-link-magnetic-wedge}
\end{equation}
Indeed, the wedge coefficient is
\((6m_T^2+L_Tq_T)/D_T^2=D_T^{-1}\), and \(dy_1,dy_2\)
are orthonormal for the fixed Euclidean quotient metric.

For \(n\in\mathbb Z^4\), set
\begin{equation}
 j_n(u)=\exp(-2\pi i n_1u_2).
 \label{cusp-link-magnetic-factor}
\end{equation}
The identity
\[
 j_m(u+n)j_n(u)=j_{m+n}(u)e^{-2\pi i m_1n_2}=j_{m+n}(u)
\]
shows that
\begin{equation}
 \mathscr L_T=(\mathbb R^4\times\mathbb C)/\mathbb Z^4,
 \qquad n\cdot(u,z)=(u+n,j_n(u)z),
 \label{cusp-link-magnetic-line}
\end{equation}
is a Hermitian line bundle on \(F_T\).  On its pullback, the unitary
connection
\begin{equation}
 \widetilde\nabla^{\mathscr L_T}=d+2\pi i u_1\,du_2
 \label{cusp-link-magnetic-line-connection}
\end{equation}
descends.  In fact, for \(\mathcal A=2\pi i u_1du_2\),
\[
 \mathcal A(u+n)=\mathcal A(u)+2\pi i n_1du_2,
 \qquad j_n^{-1}dj_n=-2\pi i n_1du_2,
\]
which is exactly the descent identity.  Its curvature is
\begin{equation}
 F_{\mathscr L_T}=2\pi i\,du_1\wedge du_2
 =2\pi i\,\theta_T^1\wedge\theta_T^2.
 \label{cusp-link-magnetic-line-curvature}
\end{equation}

Let \(\mathscr E_T=\mathscr L_T\oplus\mathscr L_T^{-1}\).  Its transition
matrices are
\[
 G_n(u)=\begin{pmatrix}j_n(u)&0\\0&j_n(u)^{-1}\end{pmatrix}
 \in\mathrm{SU}(2),
\]
so evaluation supplies a parallel determinant trivialization.  Fix the same
normalization as in \eqref{cusp-link-su2-inner-product}, and write
\begin{equation}
 H_{\mathrm{mag}}=\begin{pmatrix}i&0\\0&-i\end{pmatrix},
 \qquad |H_{\mathrm{mag}}|^2=1.
 \label{cusp-link-magnetic-generator}
\end{equation}
The direct-sum connection has the lifted expression
\begin{equation}
 \nabla^T=d+2\pi u_1\,du_2\,H_{\mathrm{mag}},
 \qquad
 F_{\nabla^T}
 =2\pi H_{\mathrm{mag}}\theta_T^1\wedge\theta_T^2
 =\frac{2\pi}{D_T}H_{\mathrm{mag}}\,dy_1\wedge dy_2.
 \label{cusp-link-magnetic-su2-connection}
\end{equation}

\begin{proposition}[Smooth critical connection and exact action]
\label{cusp-link-magnetic-critical-action}
The connection \(\nabla^T\) is non-flat and satisfies
\(d_{\nabla^T}^*F_{\nabla^T}=0\).  For
\begin{equation}
 \operatorname{YM}_T(\nabla)
 =\frac1{2g_{\mathrm{YM}}^2}
   \int_{F_T}|F_\nabla|^2\,d\operatorname{vol}_{F_T},
 \label{cusp-link-magnetic-ym-normalization}
\end{equation}
one has
\begin{equation}
 \operatorname{YM}_T(\nabla^T)
 =\frac{2\pi^2}{g_{\mathrm{YM}}^2D_T}
 =\frac{2\pi^2}{g_{\mathrm{YM}}^2T^2}
   \bigl(1+O(T^{-1})\bigr).
 \label{cusp-link-magnetic-action}
\end{equation}
This is exactly the \(M_{\mathrm{cov}}=1\) value of
\eqref{cusp-link-log-action-value}.
\end{proposition}

\begin{proof}
The transition matrices commute with \(H_{\mathrm{mag}}\), so this diagonal
endomorphism is a global parallel section of
\(\operatorname{ad}\mathscr E_T\).  The coframe is parallel for the flat
metric.  Hence \(\nabla^TF_{\nabla^T}=0\), which implies the Yang--Mills
equation.  The curvature is nonzero because \(D_T>0\).  By
\eqref{cusp-link-magnetic-wedge}, \eqref{cusp-link-cover-metric} at
\(M_{\mathrm{cov}}=1\), and \(|H_{\mathrm{mag}}|=1\),
\[
 \int_{F_T}|F_{\nabla^T}|^2\,d\operatorname{vol}_{F_T}
 =4\pi^2D_TD_T^{-2}=\frac{4\pi^2}{D_T}.
\]
This proves the exact action.  The asymptotic follows from
\(D_T=T^2+O(T)\), which is a consequence of
\eqref{exact-standing-period-assumptions}.
\end{proof}

We next specify the background-gauge Jacobi operator, including every closed
domain used below.  Let
\[
 \mathscr H_T^{\mathrm{ad},p}
 =L^2\Omega^p(F_T;\operatorname{ad}\mathscr E_T),
 \qquad p=0,1,
\]
and complexify these real Hilbert spaces for the resolvent argument.  All
operators commute with the real structure.  Define the maximal distributional
codifferential by
\begin{align}
 \operatorname{Dom}(\delta_T^{\max})
 =\{a\in\mathscr H_T^{\mathrm{ad},1}:{}&\text{ there is }u\in
 \mathscr H_T^{\mathrm{ad},0}\text{ such that}\\[-1mm]
 &\langle a,d_{\nabla^T}\phi\rangle_{L^2}
 =\langle u,\phi\rangle_{L^2}
 \text{ for all }\phi\in
 C^\infty(F_T;\operatorname{ad}\mathscr E_T)\},
 \label{cusp-link-magnetic-max-codiff-domain}\\
 \delta_T^{\max}a={}&u.\notag
\end{align}
Density of the test space makes \(u\) unique.  Passing to the limit in the
testing identity proves that \(\delta_T^{\max}\) is closed.  Therefore
\begin{equation}
 \mathscr S_T=\ker\delta_T^{\max}
 \label{cusp-link-magnetic-coulomb-slice}
\end{equation}
is a closed background Coulomb slice.

The closed scalar form
\[
 q_T^{(0)}[\phi,\psi]
 =\langle d_{\nabla^T}\phi,d_{\nabla^T}\psi\rangle_{L^2},
 \qquad
 \operatorname{Dom}(q_T^{(0)})
 =H^1\Omega^0(F_T;\operatorname{ad}\mathscr E_T),
\]
has the nonnegative Friedrichs operator
\begin{equation}
 \Delta_{\nabla^T}^{(0)}=d_{\nabla^T}^*d_{\nabla^T},
 \qquad
 \operatorname{Dom}(\Delta_{\nabla^T}^{(0)})
 =H^2\Omega^0(F_T;\operatorname{ad}\mathscr E_T).
 \label{cusp-link-magnetic-scalar-domain}
\end{equation}
The \(H^2\) equality is elliptic regularity on the compact flat torus.

In a constant oriented orthonormal frame \((e_\mu)_{\mu=1}^4\), write
\(F_{\nu\mu}=F_{\nabla^T}(e_\nu,e_\mu)\), and define
\[
 (V_Ta)_\nu=-2\sum_{\mu=1}^4[F_{\nu\mu},a_\mu].
\]
Antisymmetry of \(F_{\mu\nu}\) and skew-adjointness of
\(\operatorname{ad}(F_{\nu\mu})\) show that \(V_T\) is symmetric.  It is a
bounded bundle endomorphism.  Hence
\begin{align}
 \mathfrak j_T^{\mathrm{gf}}[a,b]
 &=\langle\nabla^Ta,\nabla^Tb\rangle_{L^2}
   +\langle V_Ta,b\rangle_{L^2},
 \label{cusp-link-magnetic-jacobi-form}\\
 \operatorname{Dom}(\mathfrak j_T^{\mathrm{gf}})
 &=H^1\Omega^1(F_T;\operatorname{ad}\mathscr E_T)\notag
\end{align}
is closed and lower bounded.  Its self-adjoint operator
\(\mathcal J_T^{\mathrm{gf}}\) has
\begin{align}
 \operatorname{Dom}(\mathcal J_T^{\mathrm{gf}})
 &=H^2\Omega^1(F_T;\operatorname{ad}\mathscr E_T),
 \label{cusp-link-magnetic-full-jacobi-domain}\\
 (\mathcal J_T^{\mathrm{gf}}a)_\nu
 &=-\sum_{\mu=1}^4\nabla_{e_\mu}\nabla_{e_\mu}a_\nu
   -2\sum_{\mu=1}^4[F_{\nu\mu},a_\mu].
 \label{cusp-link-magnetic-jacobi-expression}
\end{align}

\begin{proposition}[The Coulomb slice is reducing]
\label{cusp-link-magnetic-reduction}
For every \(z\in\mathbb C\setminus\mathbb R\) and
\(f\in\operatorname{Dom}(\delta_T^{\max})\),
\begin{equation}
 \delta_T^{\max}(\mathcal J_T^{\mathrm{gf}}-z)^{-1}f
 =(\Delta_{\nabla^T}^{(0)}-z)^{-1}\delta_T^{\max}f.
 \label{cusp-link-magnetic-resolvent-intertwining}
\end{equation}
Thus \(\mathscr S_T\) reduces \(\mathcal J_T^{\mathrm{gf}}\), and its
self-adjoint part \(\mathcal J_T\) has the exact operator and form domains
\begin{equation}
 \begin{aligned}
 \mathcal J_T&=\mathcal J_T^{\mathrm{gf}}|_{\mathscr S_T},\\
 \operatorname{Dom}(\mathcal J_T)
 &=H^2\Omega^1(F_T;\operatorname{ad}\mathscr E_T)\cap\mathscr S_T,\\
 \operatorname{Dom}(\mathfrak j_T)
 &=H^1\Omega^1(F_T;\operatorname{ad}\mathscr E_T)\cap\mathscr S_T.
 \end{aligned}
 \label{cusp-link-magnetic-reduced-domains}
\end{equation}
\end{proposition}

\begin{proof}
Write \(\nabla_\mu=\nabla^T_{e_\mu}\).  For a smooth adjoint-valued scalar
\(\phi\), the commutator
\([\nabla_\mu,\nabla_\nu]\phi=[F_{\mu\nu},\phi]\), with both curvature
terms retained, gives
\begin{align*}
 &\bigl(\mathcal J_T^{\mathrm{gf}}d_{\nabla^T}\phi
       -d_{\nabla^T}\Delta_{\nabla^T}^{(0)}\phi\bigr)_\nu\\
 &\quad=-\nabla_\mu\nabla_\mu\nabla_\nu\phi
        +\nabla_\nu\nabla_\mu\nabla_\mu\phi
        -2[F_{\nu\mu},\nabla_\mu\phi]\\
 &\quad=-[\nabla_\mu F_{\mu\nu},\phi]=0,
\end{align*}
where the last equality is the Yang--Mills equation from
Proposition~\ref{cusp-link-magnetic-critical-action}.  Formal adjunction gives
\[
 d_{\nabla^T}^*\mathcal J_T^{\mathrm{gf}}
 =\Delta_{\nabla^T}^{(0)}d_{\nabla^T}^*
\]
on smooth one-forms.

Fix \(z\notin\mathbb R\), let
\(f\in\operatorname{Dom}(\delta_T^{\max})\), and put
\(u=(\mathcal J_T^{\mathrm{gf}}-z)^{-1}f\).  The full operator domain gives
\(u\in H^2\), so \(\delta_T^{\max}u\in H^1\).  The distributional extension
of the adjoint identity gives
\[
 (\Delta_{\nabla^T}^{(0)}-z)\delta_T^{\max}u
 =\delta_T^{\max}f.
\]
The right side lies in \(L^2\); elliptic regularity upgrades
\(\delta_T^{\max}u\) to \(H^2\).  Applying the scalar resolvent proves
\eqref{cusp-link-magnetic-resolvent-intertwining}.  It follows that the
resolvents at \(z\) and \(\overline z\) preserve \(\mathscr S_T\).  By
self-adjointness they also preserve its orthogonal complement, which is the
reducing property.  The part of a self-adjoint operator on a reducing
subspace is self-adjoint, and intersection with the full operator domain gives
the second line of \eqref{cusp-link-magnetic-reduced-domains}.  The reducing
projection commutes with the form functional calculus, giving the stated form
domain.
\end{proof}

\begin{theorem}[Exact charged Landau negative mode]
\label{cusp-link-magnetic-negative-mode}
There is a smooth real \(a\in\mathscr S_T\) satisfying
\begin{equation}
 \|a\|_{L^2}=1,
 \qquad
 \mathcal J_Ta=-\frac{4\pi}{D_T}a,
 \qquad
 \operatorname{YM}_T''|_{\nabla^T}(a,a)
 =-\frac{4\pi}{g_{\mathrm{YM}}^2D_T}.
 \label{cusp-link-magnetic-negative-eigenvalue}
\end{equation}
In particular, the non-flat smooth critical connection is a classical saddle
for every finite \(T\), not a gapped local vacuum.
\end{theorem}

\begin{proof}
On the charge-two root line \(\mathscr L_T^2\), put
\[
 \mathcal D_1=\partial_{u_1},
 \qquad
 \mathcal D_2=\partial_{u_2}+4\pi i u_1,
\]
and define
\begin{equation}
 \zeta_T=\frac{q_T-6im_T}{m_T+iL_T},
 \qquad
 v_T=-\operatorname{Im}\zeta_T
 =\frac{D_T}{m_T^2+L_T^2}>0.
 \label{cusp-link-magnetic-theta-parameter}
\end{equation}
The sign is exact because
\[
 \operatorname{Im}\bigl((q_T-6im_T)(m_T-iL_T)\bigr)=-D_T.
\]
The theta series
\begin{equation}
 \Theta_T(u_1,u_2)
 =\sum_{\ell\in\mathbb Z}
 \exp\!\left(-2\pi i\zeta_T(\ell+u_1)^2+4\pi i\ell u_2\right)
 \label{cusp-link-magnetic-theta-series}
\end{equation}
converges normally with all derivatives: the absolute value of its
\(\ell\)-th term is \(e^{-2\pi v_T(\ell+u_1)^2}\).  Reindexing gives
\begin{equation}
 \Theta_T(u_1+1,u_2)=e^{-4\pi iu_2}\Theta_T(u_1,u_2),
 \qquad
 \Theta_T(u_1,u_2+1)=\Theta_T(u_1,u_2).
 \label{cusp-link-magnetic-theta-automorphy}
\end{equation}
These are precisely the factors \(j_n^2\), so \(\Theta_T\) is a smooth
section of \(\mathscr L_T^2\), independent of \(u_3,u_4\).  Termwise
differentiation gives
\begin{equation}
 (\mathcal D_1+\zeta_T\mathcal D_2)\Theta_T=0.
 \label{cusp-link-magnetic-marked-annihilator}
\end{equation}

The inverse transpose of the exact \((y_1,y_2)\)-period block in
\eqref{cusp-link-block-form} gives the physical covariant derivatives
\begin{equation}
 \nabla_{y_1}^+
 =\frac{m_T\mathcal D_1+q_T\mathcal D_2}{D_T},
 \qquad
 \nabla_{y_2}^+
 =\frac{-L_T\mathcal D_1+6m_T\mathcal D_2}{D_T}.
 \label{cusp-link-magnetic-physical-derivatives}
\end{equation}
Therefore
\begin{equation}
 \nabla_{y_1}^+-i\nabla_{y_2}^+
 =\frac{m_T+iL_T}{D_T}
  (\mathcal D_1+\zeta_T\mathcal D_2),
 \label{cusp-link-magnetic-physical-annihilator}
\end{equation}
and it kills \(\Theta_T\).  Orthogonality of the \(u_2\)-Fourier modes gives
\[
 \int_0^1\!\int_0^1|\Theta_T|^2\,du_1du_2
 =\sum_{\ell\in\mathbb Z}\int_0^1
   e^{-4\pi v_T(\ell+u_1)^2}\,du_1
 =\frac1{2\sqrt{v_T}}.
\]
Set
\begin{equation}
 \phi_T=\left(\frac{2\sqrt{v_T}}{D_T}\right)^{1/2}\Theta_T.
 \label{cusp-link-magnetic-normalized-theta}
\end{equation}
Since the marked-coordinate volume element is \(D_Td^4u\), this section has
unit \(L^2(F_T)\)-norm.  Its covariant derivatives in the \(x_1,x_2\)
directions vanish.  Furthermore,
\([\mathcal D_1,\mathcal D_2]=4\pi i\) and
\eqref{cusp-link-magnetic-physical-derivatives} imply
\[
 [\nabla_{y_1}^+,\nabla_{y_2}^+]=\frac{4\pi i}{D_T}.
\]
Writing \(\kappa_T=4\pi/D_T\), direct multiplication yields
\[
 (\nabla_{y_1}^+-i\nabla_{y_2}^+)^*
 (\nabla_{y_1}^+-i\nabla_{y_2}^+)
 =-\bigl((\nabla_{y_1}^+)^2+(\nabla_{y_2}^+)^2\bigr)-\kappa_T.
\]
Hence
\begin{equation}
 (\nabla_{y_1}^+-i\nabla_{y_2}^+)\phi_T=0,
 \qquad
 -\bigl((\nabla_{y_1}^+)^2+(\nabla_{y_2}^+)^2\bigr)\phi_T
 =\kappa_T\phi_T.
 \label{cusp-link-magnetic-ground-equation}
\end{equation}

In \(\mathfrak{sl}_2(\mathbb C)\), set
\[
 X_+=\begin{pmatrix}0&1\\0&0\end{pmatrix},
 \qquad
 X_-=\begin{pmatrix}0&0\\1&0\end{pmatrix}.
\]
Then \([H_{\mathrm{mag}},X_+]=2iX_+\), and the \(X_+\)-root line of
\(\operatorname{ad}\mathscr E_T\otimes\mathbb C\) is
\(\mathscr L_T^2\).  Thus
\[
 a^+=\phi_TX_+(dy_1-i\,dy_2)
\]
is a global complexified adjoint-valued one-form.  The first equation in
\eqref{cusp-link-magnetic-ground-equation} gives
\(d_{\nabla^T}^*a^+=0\).

Put \(\mathcal B_T=2\pi/D_T\), so that
\(\kappa_T=2\mathcal B_T\) and
\(F_{y_1y_2}=\mathcal B_TH_{\mathrm{mag}}\).  The two nonzero components are
\(a_{y_1}^+=\phi_TX_+\) and \(a_{y_2}^+=-i\phi_TX_+\).  Substitution in
\eqref{cusp-link-magnetic-jacobi-expression}, retaining the root charge two,
gives
\begin{align*}
 (\mathcal J_Ta^+)_{y_1}
 &=\kappa_T\phi_TX_+
   -2[\mathcal B_TH_{\mathrm{mag}},-i\phi_TX_+]
  =-\kappa_T\phi_TX_+,\\
 (\mathcal J_Ta^+)_{y_2}
 &=-i\kappa_T\phi_TX_+
   -2[-\mathcal B_TH_{\mathrm{mag}},\phi_TX_+]
  =i\kappa_T\phi_TX_+.
\end{align*}
Thus \(\mathcal J_Ta^+=-\kappa_Ta^+\).  For completeness, the coefficient
\(-2\) in \eqref{cusp-link-magnetic-jacobi-expression} follows by expanding
\(F_{\nabla^T+ta}=F_{\nabla^T}+t\,d_{\nabla^T}a+
\tfrac12t^2[a\wedge a]\), integrating the quadratic term by parts on the
background slice, and using the flat-base Weitzenbock identity: the Hodge
Laplacian and the curvature part of the second variation each contribute
\(-[F_{\nu\mu},a_\mu]\).

Finally take the real anti-Hermitian form
\begin{equation}
 a=\frac1{\sqrt2}\left(
   \phi_TX_+(dy_1-i\,dy_2)
   -\overline{\phi_T}X_-(dy_1+i\,dy_2)\right).
 \label{cusp-link-magnetic-real-mode}
\end{equation}
For \(X(z)=zX_+-\overline zX_-\), the fixed inner product gives
\(-\tfrac12\operatorname{Tr}(X(z)^2)=|z|^2\).  Therefore
\(|a|^2=|\phi_T|^2\), so \(\|a\|_{L^2}=1\).  Hermitian adjunction preserves
the Coulomb equation and the real Jacobi operator; hence
\(\mathcal J_Ta=-\kappa_Ta\).  On the background slice,
\[
 \operatorname{YM}_T''|_{\nabla^T}(a,a)
 =g_{\mathrm{YM}}^{-2}\langle a,\mathcal J_Ta\rangle_{L^2},
\]
which proves every identity in
\eqref{cusp-link-magnetic-negative-eigenvalue}.
\end{proof}

\subsection{Finite Gaussian quantization with independent cover and mesh}

Let
\begin{equation}
 \vartheta^j_{T,M_{\mathrm{cov}}}
 =P_{T,M_{\mathrm{cov}}}^{-\mathsf T}\mathbf e_j
 \in(\mathbb R^4)^*,
 \qquad 1\leq j\leq4.
 \label{cusp-link-physical-coframe}
\end{equation}
For \(0\leq p\leq4\), define the real finite cochain spaces
\begin{equation}
 C^p_{N_{\mathrm{mesh}}}(T,M_{\mathrm{cov}})
 =\operatorname{Map}\left(
 V_{N_{\mathrm{mesh}}},
 \bigwedge^p(\mathbb R^4)^*\otimes\mathfrak{su}(2)
 \right).
 \label{cusp-link-cochain-spaces}
\end{equation}
The exterior powers carry the metric induced by the physical coframe
\eqref{cusp-link-physical-coframe}; explicitly, for increasing multi-indices
\(I=(i_1,\ldots,i_p)\) and \(J=(j_1,\ldots,j_p)\),
\begin{equation}
 \left\langle
  \vartheta^{i_1}\wedge\cdots\wedge\vartheta^{i_p},
  \vartheta^{j_1}\wedge\cdots\wedge\vartheta^{j_p}
 \right\rangle
 =
 \det\left((G_{T,M_{\mathrm{cov}}}^{-1})_{i_a j_b}\right)_{a,b=1}^p.
 \label{cusp-link-exterior-metric}
\end{equation}
Together with \eqref{cusp-link-su2-inner-product}, this defines the
pointwise tensor-product metric.  The cochain inner product is
\begin{equation}
 \langle a,b\rangle_p
 =\frac{\mathcal D_{T,M_{\mathrm{cov}}}}{N_{\mathrm{mesh}}^4}
   \sum_{n\in V_{N_{\mathrm{mesh}}}}
   \langle a(n),b(n)\rangle.
 \label{cusp-link-cochain-inner-product}
\end{equation}
The symbol \((\vartheta^j)^\sharp\) denotes metric duality for this physical
cotangent metric, and \(\varepsilon(\vartheta^j)\) denotes left exterior
multiplication by \(\vartheta^j\).  Let
\(\mathsf T_j a(n)=a(n+\mathbf e_j)\), set
\begin{equation}
 D^+_{N_{\mathrm{mesh}},j}
 =N_{\mathrm{mesh}}(\mathsf T_j-1),
 \qquad
 D^-_{N_{\mathrm{mesh}},j}a(n)
 =N_{\mathrm{mesh}}(a(n)-a(n-\mathbf e_j)),
 \label{cusp-link-differences}
\end{equation}
and define
\begin{equation}
 d_{N_{\mathrm{mesh}}}a
 =\sum_{j=1}^4
  \varepsilon(\vartheta^j_{T,M_{\mathrm{cov}}})
  D^+_{N_{\mathrm{mesh}},j}a.
 \label{cusp-link-finite-exterior-derivative}
\end{equation}
The commuting differences and anticommuting exterior multiplications give
\(d_{N_{\mathrm{mesh}}}^2=0\).  Its adjoint is
\begin{equation}
 d_{N_{\mathrm{mesh}}}^*
 =-\sum_{j=1}^4
  \iota_{(\vartheta^j_{T,M_{\mathrm{cov}}})^\sharp}
  D^-_{N_{\mathrm{mesh}},j}.
 \label{cusp-link-finite-adjoint}
\end{equation}
Define the real reduced Coulomb space
\begin{align}
 \mathscr C_{T,M_{\mathrm{cov}},N_{\mathrm{mesh}}}
 &=
 \ker(d_{N_{\mathrm{mesh}}}^*:C^1\to C^0)
 \cap
 \left(\ker\Delta^{(1)}_{N_{\mathrm{mesh}}}\right)^\perp,
 \label{cusp-link-reduced-coulomb-space}\\
 \Delta^{(1)}_{N_{\mathrm{mesh}}}
 &=d_{N_{\mathrm{mesh}}}d_{N_{\mathrm{mesh}}}^*
 +d_{N_{\mathrm{mesh}}}^*d_{N_{\mathrm{mesh}}}.
 \label{cusp-link-finite-hodge-laplacian}
\end{align}
On this finite-dimensional space put
\begin{equation}
 \mathcal L_{T,M_{\mathrm{cov}},N_{\mathrm{mesh}}}
 =g_{\mathrm{YM}}^{-2}d_{N_{\mathrm{mesh}}}^*
                         d_{N_{\mathrm{mesh}}}.
 \label{cusp-link-finite-hessian}
\end{equation}

\begin{proposition}[Exact finite Hessian spectrum]
\label{cusp-link-finite-spectrum}
For
\begin{equation}
 \mathfrak d_{N_{\mathrm{mesh}}}(k)_j
 =N_{\mathrm{mesh}}
  \left(e^{2\pi i k_j/N_{\mathrm{mesh}}}-1\right),
 \qquad
 k\in(\mathbb Z/N_{\mathrm{mesh}}\mathbb Z)^4,
 \label{cusp-link-fourier-symbol}
\end{equation}
In the Fourier formulas, \(\lVert\,\cdot\,\rVert\) on the complexified physical
cotangent space denotes the Hermitian extension of the Euclidean norm.  The
positive spectrum of \(\mathcal L\) is obtained from
\begin{equation}
 \lambda_{T,M_{\mathrm{cov}},N_{\mathrm{mesh}}}(k)
 =\frac1{g_{\mathrm{YM}}^2}
  \left\lVert
   P_T^{-\mathsf T}S_{M_{\mathrm{cov}}}^{-1}
   \mathfrak d_{N_{\mathrm{mesh}}}(k)
  \right\rVert^2,
 \qquad k\neq0.
 \label{cusp-link-finite-mode-eigenvalue}
\end{equation}
Each complex Fourier fibre has transverse multiplicity
\(3\dim\mathfrak{su}(2)=9\).  In particular,
\begin{equation}
 \lambda_{\min}(\mathcal L)
 =\min_{k\neq0}\lambda_{T,M_{\mathrm{cov}},N_{\mathrm{mesh}}}(k)>0.
 \label{cusp-link-finite-minimum}
\end{equation}
For \(\mathbf e_1=(1,0,0,0)^{\mathsf T}\), the exact mode value is
\begin{equation}
 \lambda_{T,M_{\mathrm{cov}},N_{\mathrm{mesh}}}(\mathbf e_1)
 =\frac{4N_{\mathrm{mesh}}^2
              \sin^2(\pi/N_{\mathrm{mesh}})}{g_{\mathrm{YM}}^2}
  \frac{L_T^2+m_T^2}{D_T^2};
 \label{cusp-link-cover-soft-mode}
\end{equation}
it is independent of \(M_{\mathrm{cov}}\).
\end{proposition}

\begin{proof}
Translation is unitary for \eqref{cusp-link-cochain-inner-product}, so
\((D_j^+)^*=-D_j^-\).  Exterior multiplication by \(\vartheta^j\) has
adjoint contraction by \((\vartheta^j)^\sharp\), proving
\eqref{cusp-link-finite-adjoint}.  On the Fourier character
\(\chi_k(n)=e^{2\pi i k\cdot n/N_{\mathrm{mesh}}}\), the derivative symbol
is exterior multiplication by
\[
 \zeta(k)=P_{T,M_{\mathrm{cov}}}^{-\mathsf T}
           \mathfrak d_{N_{\mathrm{mesh}}}(k)
 =P_T^{-\mathsf T}S_{M_{\mathrm{cov}}}^{-1}
           \mathfrak d_{N_{\mathrm{mesh}}}(k).
\]
The Hermitian exterior-algebra identity
\[
 \iota_{\overline\zeta^\sharp}\varepsilon(\zeta)
 +\varepsilon(\zeta)\iota_{\overline\zeta^\sharp}
 =\lVert\zeta\rVert^2\operatorname{id}
\]
shows that on the Coulomb polarization
\(\ker\iota_{\overline\zeta^\sharp}\subset(\mathbb C^4)^*\), the operator
\(d^*d\) is multiplication by \(\lVert\zeta\rVert^2\).  For \(k\neq0\),
\(\mathfrak d_{N_{\mathrm{mesh}}}(k)\neq0\), so invertibility of
\(P_{T,M_{\mathrm{cov}}}\) gives
\(\zeta(k)\neq0\).  The contraction has rank one, and its kernel has
dimension three.  The zero Fourier fibre consists exactly of constant
harmonic cochains and was removed in
\eqref{cusp-link-reduced-coulomb-space}.  Fourier completeness now proves
\eqref{cusp-link-finite-mode-eigenvalue}--\eqref{cusp-link-finite-minimum}.

For \(k=\mathbf e_1\), the only nonzero component of
\(\mathfrak d_{N_{\mathrm{mesh}}}(k)\) has modulus
\[
 2N_{\mathrm{mesh}}\sin(\pi/N_{\mathrm{mesh}}),
\]
and
\(S_{M_{\mathrm{cov}}}^{-1}\mathbf e_1=\mathbf e_1\).  Solving
\(P_T^{\mathsf T}\xi=\mathbf e_1\), or directly inverting the lower
\(2\times2\) block in \eqref{cusp-link-period-matrix}, gives
\[
 \lVert P_T^{-\mathsf T}\mathbf e_1\rVert^2
 =\frac{L_T^2+m_T^2}{D_T^2}.
\]
Substitution proves \eqref{cusp-link-cover-soft-mode}.
\end{proof}

\subsection{Native continuum comparison and the exact all-mesh gap}

The cover formula \eqref{cusp-link-finite-minimum} is already exact for every
\(M_{\mathrm{cov}}\).  On the native torus, meaning
\(M_{\mathrm{cov}}=1\), its minimum can be evaluated completely and compared
with the continuum Yang--Mills Hessian.  Let \(A_0\) be the product connection
on \(F_T\times\mathrm{SU}(2)\), and normalize
\begin{equation}
 \operatorname{YM}_T(A)
 =\frac1{2g_{\mathrm{YM}}^2}
   \int_{F_T}|F_A|^2\,\mathrm d\operatorname{vol}_{F_T}.
 \label{cusp-link-native-ym}
\end{equation}
On one-forms let \(\mathrm d_1\) be the maximal closed exterior derivative
\[
 \mathrm d_1:L^2\Omega^1(F_T;\mathfrak{su}(2))
 \longrightarrow L^2\Omega^2(F_T;\mathfrak{su}(2)),
\]
with domain
\begin{equation}
 \operatorname{Dom}(\mathrm d_1)
 =\{a\in L^2\Omega^1:\mathrm da\in L^2\Omega^2
       \text{ distributionally}\}.
 \label{cusp-link-native-maximal-d}
\end{equation}

\begin{proposition}[Exact native continuum Hessian and Coulomb spectrum]
\label{cusp-link-native-continuum-hessian}
For smooth one-forms \(a,b\),
\begin{equation}
 \operatorname{Hess}_{A_0}\operatorname{YM}_T(a,b)
 =g_{\mathrm{YM}}^{-2}\langle\mathrm da,\mathrm db\rangle_{L^2}.
 \label{cusp-link-native-hessian-formula}
\end{equation}
The associated closed form and nonnegative self-adjoint operator are
\begin{align}
 \mathfrak l_T[a]&=g_{\mathrm{YM}}^{-2}\|\mathrm d_1a\|_{L^2}^2,
 &\operatorname{Dom}(\mathfrak l_T)&=\operatorname{Dom}(\mathrm d_1),
 \label{cusp-link-native-closed-form}\\
 \mathcal L_T^{\mathrm{cont}}
 &=g_{\mathrm{YM}}^{-2}\mathrm d_1^*\mathrm d_1,
 &\operatorname{Dom}(\mathcal L_T^{\mathrm{cont}})
 &=\{a\in\operatorname{Dom}(\mathrm d_1):
      \mathrm d_1a\in\operatorname{Dom}(\mathrm d_1^*)\}.
 \label{cusp-link-native-continuum-operator}
\end{align}
The reduced Coulomb space is the closed subspace
\begin{equation}
 \mathscr C_T^{\mathrm{cont}}
 =\{a\in L^2\Omega^1(F_T;\mathfrak{su}(2)):
      \mathrm d^*a=0\ \text{distributionally},\quad
      a\perp\mathcal H^1(F_T;\mathfrak{su}(2))\}.
 \label{cusp-link-native-continuum-coulomb}
\end{equation}
It reduces \(\mathcal L_T^{\mathrm{cont}}\).  The restricted form domain is
\(H^1\Omega^1\cap\mathscr C_T^{\mathrm{cont}}\), and the restricted operator
domain is \(H^2\Omega^1\cap\mathscr C_T^{\mathrm{cont}}\).

For \(k\in\mathbb Z^4\setminus\{0\}\), the character
\(\chi_{T,k}([x])=\exp(2\pi i\langle P_T^{-\mathsf T}k,x\rangle)\)
has Coulomb polarization
\begin{equation}
 \ker\iota_{(P_T^{-\mathsf T}k)^\sharp}
 \otimes\mathfrak{su}(2)_{\mathbb C}
 \label{cusp-link-native-continuum-polarization}
\end{equation}
of complex dimension nine, and the operator eigenvalue on that fibre is
\begin{equation}
 \lambda_T^{\mathrm{cont}}(k)
 =\frac{4\pi^2}{g_{\mathrm{YM}}^2}|P_T^{-\mathsf T}k|^2.
 \label{cusp-link-native-continuum-eigenvalue}
\end{equation}
For all sufficiently large \(T\), put
\begin{equation}
 a_{\gamma,T}=\frac{L_T^2+m_T^2}{D_T^2},
 \qquad
 a_{u,T}=\frac{q_T^2+36m_T^2}{D_T^2},
 \qquad
 I_T=\{j\in\{\gamma,u\}:a_{j,T}=\min(a_{\gamma,T},a_{u,T})\}.
 \label{cusp-link-native-short-data}
\end{equation}
Then the first reduced continuum eigenvalue and its real multiplicity are
\begin{equation}
 \gamma_T^{\mathrm{cont}}
 =\frac{4\pi^2}{g_{\mathrm{YM}}^2}
   \min\{a_{\gamma,T},a_{u,T}\},
 \qquad
 \operatorname{mult}_{\mathbb R}(\gamma_T^{\mathrm{cont}})=18\,\#I_T.
 \label{cusp-link-native-continuum-gap}
\end{equation}
\end{proposition}

\begin{proof}
With \(F_{A_0+a}=\mathrm da+\tfrac12[a\wedge a]\), the coefficient of
\(st\) in \(|F_{A_0+sa+tb}|^2\) is
\(2\langle\mathrm da,\mathrm db\rangle\), which proves
\eqref{cusp-link-native-hessian-formula}.  The maximal distributional
derivative is closed and densely defined.  The first representation theorem
for its graph-norm square therefore gives
\eqref{cusp-link-native-closed-form}--\eqref{cusp-link-native-continuum-operator}
with exactly the displayed domains.

On the flat torus the Fourier identities are
\begin{align*}
 \mathrm d(\chi_{T,k}\alpha)
 &=2\pi i\chi_{T,k}\varepsilon(\xi_{T,k})\alpha,\\
 \mathrm d^*(\chi_{T,k}\alpha)
 &=-2\pi i\chi_{T,k}\iota_{\xi_{T,k}^\sharp}\alpha,\\
 \mathrm d^*\mathrm d(\chi_{T,k}\alpha)
 &=4\pi^2\chi_{T,k}
   \bigl(|\xi_{T,k}|^2\alpha
   -\varepsilon(\xi_{T,k})\iota_{\xi_{T,k}^\sharp}\alpha\bigr),
 \qquad \xi_{T,k}=P_T^{-\mathsf T}k.
\end{align*}
They show directly that \(\mathscr C_T^{\mathrm{cont}}\) reduces the operator.
They also identify its form and operator domains respectively with
\(H^1\cap\mathscr C_T^{\mathrm{cont}}\) and
\(H^2\cap\mathscr C_T^{\mathrm{cont}}\).  For \(k\ne0\), contraction by
\(\xi_{T,k}^\sharp\) has rank one on the four-dimensional complexified
cotangent space.  Its kernel therefore has dimension three; tensoring with
\(\mathfrak{su}(2)_{\mathbb C}\) gives dimension nine, and the last displayed
identity proves \eqref{cusp-link-native-continuum-eigenvalue}.  The zero mode
is precisely the constant harmonic sector removed in
\eqref{cusp-link-native-continuum-coulomb}.

It remains to identify the shortest nonzero dual vectors.  Write
\(k=(r,s,n_1,n_2)^{\mathsf T}\).  Solving \(P_T^{\mathsf T}x=k\), the third
and fourth equations give \(x_1=n_1\), \(x_3=n_2\).  Thus
\(|P_T^{-\mathsf T}k|\ge1\) when \((n_1,n_2)\ne(0,0)\).  When those two
indices vanish,
\begin{equation}
 P_T^{-\mathsf T}(r,s,0,0)^{\mathsf T}
 =\left(0,\frac{m_Tr+q_Ts}{D_T},0,
           \frac{-L_Tr+6m_Ts}{D_T}\right).
 \label{cusp-link-native-soft-plane}
\end{equation}
The scaled two-dimensional matrix in this expression satisfies
\[
 \frac{T}{D_T}
 \begin{pmatrix}m_T&q_T\\-L_T&6m_T\end{pmatrix}
 =\begin{pmatrix}0&1\\-1&0\end{pmatrix}+O_{\mathrm{op}}(T^{-1}).
\]
Hence, uniformly in real \((r,s)\),
\[
 (1-C/T)\sqrt{r^2+s^2}
 \le T|P_T^{-\mathsf T}(r,s,0,0)^{\mathsf T}|
 \le(1+C/T)\sqrt{r^2+s^2}.
\]
For large \(T\), every integer pair of squared norm at least two is longer
than each unit pair, while every mode with \((n_1,n_2)\ne0\) has length at
least one.  The only possible shortest vectors are therefore the two signed
pairs with indices \(\pm(1,0,0,0)\) and \(\pm(0,1,0,0)\).  Their squared
lengths, read from \eqref{cusp-link-native-soft-plane}, are precisely
\(a_{\gamma,T}\) and \(a_{u,T}\).  Each signed pair contributes eighteen
real dimensions, proving \eqref{cusp-link-native-continuum-gap}.
\end{proof}

Set
\begin{equation}
 \rho_N=4N^2\sin^2(\pi/N),\qquad
 \eta_N=\begin{cases}1,&N=2,\\2,&N>2.\end{cases}
 \label{cusp-link-native-rho-eta}
\end{equation}

\begin{theorem}[Exact native finite gap in every mesh regime]
\label{cusp-link-native-finite-gap}
There is \(T_1\), independent of \(N\ge2\), such that for \(T\ge T_1\),
the first eigenvalue of
\(\mathcal L_{T,1,N}\) on \(\mathscr C_{T,1,N}\) is
\begin{equation}
 \gamma_{N,T}^{\mathrm{fin}}
 =\frac{\rho_N}{g_{\mathrm{YM}}^2}
   \min\{a_{\gamma,T},a_{u,T}\}
 =\frac{\rho_N}{4\pi^2}\gamma_T^{\mathrm{cont}}.
 \label{cusp-link-native-finite-gap-exact}
\end{equation}
Its real multiplicity is exactly
\begin{equation}
 \operatorname{mult}_{\mathbb R}(\gamma_{N,T}^{\mathrm{fin}})
 =9\eta_N\,\#I_T.
 \label{cusp-link-native-finite-gap-multiplicity}
\end{equation}
Uniformly for every integer-valued choice \(N(T)\ge2\), there are constants
\(0<c<C<\infty\) such that
\begin{equation}
 \frac{c}{T^2}\le\gamma_{N(T),T}^{\mathrm{fin}}
 \le\frac{C}{T^2}
 \label{cusp-link-native-finite-gap-uniform}
\end{equation}
for all large \(T\).  If \(N(T)\to\infty\), then
\begin{equation}
 \gamma_{N(T),T}^{\mathrm{fin}}
 =\frac{4\pi^2}{g_{\mathrm{YM}}^2T^2}
   \left(1+O(T^{-1})+O(N(T)^{-2})\right).
 \label{cusp-link-native-finite-gap-refined}
\end{equation}
The largest physical edge in the native cubulation is
\begin{equation}
 \max_{1\le j\le4}\frac{|P_T\mathbf e_j|}{N(T)}
 =\frac{T}{N(T)}(1+O(T^{-1})),
 \label{cusp-link-native-mesh-size}
\end{equation}
so this varying continuum torus is resolved exactly when
\(N(T)/T\to\infty\).  This geometric condition is not needed for the
algebraic gap bounds above.
\end{theorem}

\begin{proof}
Put \(\sigma_N=2N\sin(\pi/N)\), so \(\rho_N=\sigma_N^2\).  The least
nonzero modulus of a component of \(\mathfrak d_N(k)\) is \(\sigma_N\),
attained exactly at frequency \(\pm1\), with the two signs equal when
\(N=2\).  If \(\mathfrak d_N(k)\) is supported in its first two components
but \(k\notin\{\pm\mathbf e_1,\pm\mathbf e_2\}\), then
\begin{equation}
 |\mathfrak d_N(k)_1|^2+|\mathfrak d_N(k)_2|^2\ge2\sigma_N^2.
 \label{cusp-link-native-next-frequency}
\end{equation}
Indeed, either both components are nonzero, or the single nonzero frequency
is not \(\pm1\); in the latter case
\(\sin(2\pi/N)/\sin(\pi/N)=2\cos(\pi/N)\ge\sqrt2\) for \(N\ge4\), and
\(N=2,3\) are immediate.

On the first two components, \eqref{cusp-link-native-soft-plane} and the
uniform singular-value estimate in its proof show that the factor
\(\sqrt2\) in \eqref{cusp-link-native-next-frequency} separates every other
frequency from the two coordinate frequencies, uniformly in \(N\), once
\(T\) is large.  For a general complex vector \(z\in\mathbb C^4\), the
third and fourth equations in \(P_T^{\mathsf T}x=z\) give
\(x_1=z_3\), \(x_3=z_4\).  Thus any finite frequency with nonzero third or
fourth component satisfies
\[
 |P_T^{-\mathsf T}\mathfrak d_N(k)|^2
 \ge|\mathfrak d_N(k)_3|^2+|\mathfrak d_N(k)_4|^2
 \ge\sigma_N^2,
\]
whereas either coordinate candidate is \(O(\sigma_N^2/T^2)\).  Therefore the
only minimizers are the signed coordinate pairs selected by \(I_T\).
The exact mode formula \eqref{cusp-link-finite-mode-eigenvalue} gives
\eqref{cusp-link-native-finite-gap-exact}.  A signed complex Fourier fibre has
dimension nine; reality pairs the two signs when \(N>2\), while they coincide
when \(N=2\).  This proves
\eqref{cusp-link-native-finite-gap-multiplicity}.

Finally, the standing assumptions give
\(a_{\gamma,T}=T^{-2}(1+O(T^{-1}))\) and
\(a_{u,T}=T^{-2}(1+O(T^{-1}))\).  The elementary sine bounds yield
\(16\le\rho_N\le4\pi^2\), proving
\eqref{cusp-link-native-finite-gap-uniform}.  When \(N(T)\to\infty\),
\(\rho_{N(T)}=4\pi^2(1+O(N(T)^{-2}))\), which proves
\eqref{cusp-link-native-finite-gap-refined}.  The first two columns of
\(P_T\) have norm \(T(1+O(T^{-1}))\), while its last two columns have norm
one.  Division by \(N(T)\) proves \eqref{cusp-link-native-mesh-size} and the
final assertion.
\end{proof}

\begin{remark}[Limits of the comparison]
\label{cusp-link-native-comparison-scope}
The exact minimum in Theorem~\ref{cusp-link-native-finite-gap} is a statement
only for \(M_{\mathrm{cov}}=1\).  For a nontrivial cover,
\(S_{M_{\mathrm{cov}}}^{-1}\) creates additional frequency scales, so the
general statement remains the exact minimization
\eqref{cusp-link-finite-minimum} together with proved mode bounds such as
\eqref{cusp-link-cover-soft-mode}.  Moreover,
\(\mathcal L_T^{\mathrm{cont}}\) is a classical Hessian operator; no
quantization or unitary identification with the finite Gaussian Hilbert space
is contained in Proposition~\ref{cusp-link-native-continuum-hessian}.
\end{remark}

Let \(da\) denote Lebesgue measure on
\(\mathscr C_{T,M_{\mathrm{cov}},N_{\mathrm{mesh}}}\) induced by
\eqref{cusp-link-cochain-inner-product} and define
\begin{equation}
 \mathcal H^{\mathrm{Gauss}}_{T,M_{\mathrm{cov}},N_{\mathrm{mesh}}}
 =L^2(\mathscr C_{T,M_{\mathrm{cov}},N_{\mathrm{mesh}}},da).
 \label{cusp-link-gaussian-hilbert-space}
\end{equation}
Put
\(r=\dim_{\mathbb R}\mathscr C_{T,M_{\mathrm{cov}},N_{\mathrm{mesh}}}\).
In any orthonormal linear coordinates \((x_1,\ldots,x_r)\) for the inner
product induced by \eqref{cusp-link-cochain-inner-product}, define
\begin{equation}
 \Delta_{\mathscr C}=\sum_{s=1}^r\frac{\partial^2}{\partial x_s^2}.
 \label{cusp-link-coulomb-laplacian}
\end{equation}
This definition is independent of the chosen orthonormal coordinates.

\begin{theorem}[Finite Gaussian oscillator and a resolved cusp sequence]
\label{cusp-link-gaussian-theorem}
On the initial domain
\(C_c^\infty(\mathscr C_{T,M_{\mathrm{cov}},N_{\mathrm{mesh}}};\mathbb C)\),
consider the symmetric quadratic differential operator
\begin{equation}
 \widehat H^{\mathrm{Gauss}}_{T,M_{\mathrm{cov}},N_{\mathrm{mesh}}}
 =-\frac12\Delta_{\mathscr C}
  +\frac12\langle a,
     \mathcal L_{T,M_{\mathrm{cov}},N_{\mathrm{mesh}}}a\rangle
 \label{cusp-link-gaussian-hamiltonian}
\end{equation}
where the second summand denotes multiplication by the displayed quadratic
function of \(a\).
Its Friedrichs realization on \eqref{cusp-link-gaussian-hilbert-space} has a
unique Gaussian ground state and spectral gap
\begin{equation}
 \operatorname{gap}(\widehat H^{\mathrm{Gauss}})
 =\sqrt{\lambda_{\min}(
   \mathcal L_{T,M_{\mathrm{cov}},N_{\mathrm{mesh}}})}.
 \label{cusp-link-gaussian-gap}
\end{equation}
Consequently,
\begin{equation}
 \operatorname{gap}(\widehat H^{\mathrm{Gauss}})
 \leq
 \frac{2N_{\mathrm{mesh}}\sin(\pi/N_{\mathrm{mesh}})}{g_{\mathrm{YM}}}
 \frac{\sqrt{L_T^2+m_T^2}}{D_T}.
 \label{cusp-link-gaussian-gap-bound}
\end{equation}
For \(M_{\mathrm{cov}}=1\) and all sufficiently large \(T\), the exact
native formula is
\begin{equation}
 \operatorname{gap}(\widehat H^{\mathrm{Gauss}}_{T,1,N_{\mathrm{mesh}}})
 =\frac{2N_{\mathrm{mesh}}\sin(\pi/N_{\mathrm{mesh}})}{g_{\mathrm{YM}}}
  \sqrt{\min\{a_{\gamma,T},a_{u,T}\}}.
 \label{cusp-link-native-gaussian-gap}
\end{equation}
The complex dimension of the first-excited eigenspace is
\(9\eta_{N_{\mathrm{mesh}}}\#I_T\).  Uniformly for every
\(N_{\mathrm{mesh}}(T)\ge2\), this gap is comparable to \(T^{-1}\); if
\(N_{\mathrm{mesh}}(T)\to\infty\), it equals
\begin{equation}
 \frac{2\pi}{g_{\mathrm{YM}}T}
 \left(1+O(T^{-1})+O(N_{\mathrm{mesh}}(T)^{-2})\right).
 \label{cusp-link-native-gaussian-gap-asymptotic}
\end{equation}

There is an explicit diagonal sequence for which the cover geometry tends locally to
\(\mathbb R^4\), the cubical mesh resolves that geometry, and the Gaussian
gap tends to zero.  Explicitly, for integers \(j\to\infty\), take
\begin{equation}
 T_j=j^2,\qquad M_{\mathrm{cov},j}=j,\qquad
 N_{\mathrm{mesh},j}=j^3.
 \label{cusp-link-diagonal-sequence}
\end{equation}
Then
\begin{equation}
 \begin{aligned}
 \operatorname{sys}(\widetilde F_{T_j,M_{\mathrm{cov},j}})
 &\longrightarrow\infty,\\
 \max_r
 \frac{\lVert P_{T_j,M_{\mathrm{cov},j}}\mathbf e_r\rVert}
      {N_{\mathrm{mesh},j}}
 &\longrightarrow0,\\
 \operatorname{gap}(\widehat H^{\mathrm{Gauss}})
 &\longrightarrow0,\\
 S^{\log}_{T_j,M_{\mathrm{cov},j},N_{\mathrm{mesh},j}}(U)
 &\longrightarrow0.
 \end{aligned}
 \label{cusp-link-diagonal-conclusions}
\end{equation}
Thus the pointed cover tori converge locally isometrically to \(\mathbb R^4\)
while their finite linearized Gaussian systems are simultaneously refined.
\end{theorem}

\begin{proof}
The real finite-dimensional operator \(\mathcal L\) is positive definite on
the reduced Coulomb space.  Choose an orthonormal eigenbasis with eigenvalues
\(0<\lambda_1\leq\cdots\leq\lambda_r\).  In the resulting coordinates,
\eqref{cusp-link-gaussian-hamiltonian} is the tensor sum
\[
 \sum_{s=1}^r
 \left(-\frac12\frac{\partial^2}{\partial x_s^2}
       +\frac12\lambda_sx_s^2\right).
\]
The normalized Hermite basis has eigenvalues
\(\sum_s(n_s+\tfrac12)\sqrt{\lambda_s}\), with
\(n_s\in\mathbb Z_{\geq0}\).  The all-zero multi-index is the unique ground
state, and the smallest possible positive excitation is
\(\sqrt{\lambda_1}\).  This proves \eqref{cusp-link-gaussian-gap}; applying
the minimum to the explicit \(\mathbf e_1\)-mode in
\eqref{cusp-link-cover-soft-mode} proves
\eqref{cusp-link-gaussian-gap-bound}.  When \(M_{\mathrm{cov}}=1\),
Theorem~\ref{cusp-link-native-finite-gap} evaluates \(\lambda_1\) and its
real eigenspace exactly.  Taking its positive square root proves
\eqref{cusp-link-native-gaussian-gap}; the same operation on
\eqref{cusp-link-native-finite-gap-uniform} and
\eqref{cusp-link-native-finite-gap-refined} gives the stated comparison and
\eqref{cusp-link-native-gaussian-gap-asymptotic}.

Since \(T_j,M_{\mathrm{cov},j}\to\infty\),
Proposition~\ref{cusp-link-pointed-limits} proves both the systole conclusion
and the pointed local-isometric \(\mathbb R^4\) limit.

The first two columns of \(P_T\) have norm \(O(T)\), while the last two
columns of \(P_{T,M_{\mathrm{cov}}}\) have norm exactly
\(M_{\mathrm{cov}}\).  Therefore the largest physical edge in the
\(N_{\mathrm{mesh}}\)-subdivision has length
\(O(\max\{T,M_{\mathrm{cov}}\}/N_{\mathrm{mesh}})\), which is \(O(j^{-1})\)
for \eqref{cusp-link-diagonal-sequence}.  Finally,
\eqref{cusp-link-cusp-asymptotics} gives
\(\sqrt{L_T^2+m_T^2}/D_T=T^{-1}(1+O(T^{-1}))\), while
\(2N\sin(\pi/N)\leq2\pi\).  The right-hand side of
\eqref{cusp-link-gaussian-gap-bound} is therefore \(O(j^{-2})\).
Moreover, \eqref{cusp-link-log-action-value} is \(O(j^{-2})\), since
\(M_{\mathrm{cov},j}^2/D_{T_j}=j^2/(j^4+O(j^2))\).
This proves every assertion in \eqref{cusp-link-diagonal-conclusions}.
\end{proof}

\subsection{A gapless auxiliary direct-integral operator}
\label{cusp-link-auxiliary-direct-integral}

The closing of the native scalar eigenvalues can be attached to one operator
without identifying that operator with a physical Hamiltonian.  Use the
coordinate diffeomorphism \([y]\mapsto[P_Ty]\) from
\eqref{cusp-link-coordinate-diffeomorphism}; normalized Haar measure on
\(F_T\) pulls back to normalized Lebesgue measure \(dy\) on
\(\mathbb T^4=\mathbb R^4/\mathbb Z^4\).  For
\(k\in\mathbb Z^4\setminus\{0\}\), put
\[
 e_k(y)=e^{2\pi i k\cdot y},
 \qquad
 \lambda_k(T)=4\pi^2|P_T^{-\mathsf T}k|^2.
\]
Fix \(T_1\ge T_0\) sufficiently large for the estimates following
\eqref{cusp-link-native-soft-plane}, and define
\begin{align}
 \mathscr H_{\mathrm{aux}}
 &=L^2\!\left([T_1,\infty),dT;
       L^2_0(\mathbb T^4,dy)\right)
 \cong\bigoplus_{k\in\mathbb Z^4\setminus\{0\}}
       L^2([T_1,\infty),dT),
 \label{cusp-link-auxiliary-Hilbert-space}\\
 (\mathscr L_{\mathrm{aux}}u)_k(T)
 &=\lambda_k(T)u_k(T),
 \label{cusp-link-auxiliary-operator}\\
 \operatorname{Dom}(\mathscr L_{\mathrm{aux}})
 &=\left\{u\in\mathscr H_{\mathrm{aux}}:
   \sum_{k\ne0}\int_{T_1}^{\infty}
   \lambda_k(T)^2|u_k(T)|^2\,dT<\infty\right\}.
 \label{cusp-link-auxiliary-domain}
\end{align}
Here \(L^2_0\) is the orthogonal complement of the constants.  Thus
\(\mathscr L_{\mathrm{aux}}\) is, by definition, the direct integral of the
mean-zero scalar Laplacians on the native fibres, expressed on a fixed
Hilbert space.

\begin{theorem}[Exact zero essential threshold]
\label{cusp-link-auxiliary-gapless-theorem}
The operator \(\mathscr L_{\mathrm{aux}}\) is nonnegative and self-adjoint on
the domain \eqref{cusp-link-auxiliary-domain}.  Its kernel is zero, but
\begin{equation}
 0\in\operatorname{Spec}_{\mathrm{ess}}(\mathscr L_{\mathrm{aux}}).
 \label{cusp-link-auxiliary-essential-zero}
\end{equation}
It therefore has no positive spectral threshold above zero.
\end{theorem}

\begin{proof}
For each \(k\ne0\), \(\lambda_k\) is a positive continuous real function.
Multiplication by \(\lambda_k\) on
\[
 \{v\in L^2([T_1,\infty)):\lambda_kv\in L^2([T_1,\infty))\}
\]
is nonnegative and self-adjoint: testing the adjoint against bounded compactly
supported functions forces \(\lambda_kv\in L^2\), and the adjoint acts by the
same real multiplier.  The orthogonal sum of these multiplication operators
has exactly the domain \eqref{cusp-link-auxiliary-domain}, proving
self-adjointness.  Strict positivity of every \(\lambda_k(T)\) proves that the
kernel is zero.

Take \(k_0=(1,0,0,0)\).  Formula
\eqref{cusp-link-native-soft-plane} and the standing cusp bounds give
\begin{equation}
 \lambda_{k_0}(T)
 =4\pi^2\frac{L_T^2+m_T^2}{D_T^2}
 =\frac{4\pi^2}{T^2}\bigl(1+O(T^{-1})\bigr).
 \label{cusp-link-auxiliary-soft-eigenvalue}
\end{equation}
For integers \(n>T_1\), let \(u^{(n)}\) have only its \(k_0\)-component
nonzero, with
\[
 u^{(n)}_{k_0}(T)=\mathbf1_{[n,n+1]}(T).
\]
The \(u^{(n)}\) are orthonormal members of
\(\operatorname{Dom}(\mathscr L_{\mathrm{aux}})\), converge weakly to zero,
and satisfy
\[
 \|\mathscr L_{\mathrm{aux}}u^{(n)}\|^2
 =\int_n^{n+1}\lambda_{k_0}(T)^2\,dT
 \le\sup_{T\in[n,n+1]}\lambda_{k_0}(T)^2
 \longrightarrow0.
\]
This is a Weyl singular sequence at zero.  Because the kernel is zero, it
places zero in the essential rather than discrete spectrum, proving
\eqref{cusp-link-auxiliary-essential-zero}.
\end{proof}

This theorem concerns an auxiliary scalar direct integral whose base measure
is \(dT\).  It is neither the finite Gaussian oscillator nor a Yang--Mills
Hamiltonian, and it supplies no map between their Hilbert spaces.

\subsection{The exponential map and the exact physical boundary}

Let \(C^1=C^1_{N_{\mathrm{mesh}}}(T,M_{\mathrm{cov}})\), regarded only as a
finite-dimensional Lie-algebra cochain space, and define componentwise
\begin{equation}
 \operatorname{Exp}_{N_{\mathrm{mesh}}}:C^1
 \longrightarrow\mathrm{SU}(2)^{E_{N_{\mathrm{mesh}}}},
 \qquad
 a\longmapsto\left(\exp(a_e/N_{\mathrm{mesh}})\right)_e.
 \label{cusp-link-exponential-map}
\end{equation}
Here, if
\(a(n)=\sum_{j=1}^4\vartheta^j_{T,M_{\mathrm{cov}}}\otimes a_j(n)\),
then \(a_e=a_j(n)\) for the oriented edge \(e=(n,j)\).

\begin{proposition}[Local chart, not a global embedding]
\label{cusp-link-exponential-local-only}
The map \eqref{cusp-link-exponential-map} is smooth and a local
diffeomorphism at zero, but it is not globally injective.  Under left
trivialization of every tangent space at the identity,
\begin{equation}
 D\operatorname{Exp}_{N_{\mathrm{mesh}}}(0)\,a
 =N_{\mathrm{mesh}}^{-1}a.
 \label{cusp-link-exponential-derivative}
\end{equation}
If
\(S_W:\mathrm{SU}(2)^{E_{N_{\mathrm{mesh}}}}\to\mathbb R\) is any smooth
compact-link action for which the identity configuration
\(\mathbf1=(1)_{e\in E_{N_{\mathrm{mesh}}}}\) is critical, then
\begin{equation}
 D^2(S_W\circ\operatorname{Exp}_{N_{\mathrm{mesh}}})_0[a,b]
 =N_{\mathrm{mesh}}^{-2}
  D^2(S_W)_{\mathbf1}[a,b]
 \label{cusp-link-exponential-hessian-scaling}
\end{equation}
under the same tangent identification.  Thus, for a prescribed cochain
quadratic form \(Q\),
\[
 D^2(S_W\circ\operatorname{Exp}_{N_{\mathrm{mesh}}})_0=Q
\]
holds exactly when
\[
 D^2(S_W)_{\mathbf1}=N_{\mathrm{mesh}}^2Q
\]
under the coefficient identification.  Both this scaling and the tensor
structure must be verified for the chosen compact-link action; neither
follows from the exponential map alone.
\end{proposition}

\begin{proof}
For the Lie-group exponential, \(D\exp_0\) is the identity.  Applying the
chain rule to \(a_e\mapsto a_e/N_{\mathrm{mesh}}\) gives
\eqref{cusp-link-exponential-derivative}, which is invertible, so the inverse
function theorem gives a local diffeomorphism.  It is not globally
injective: a cochain which is zero on every edge except one, where it equals
\(2\pi N_{\mathrm{mesh}}H_{\mathrm{link}}\), has the same image as the zero cochain by
\eqref{cusp-link-flux-generator}.  Finally, the second-derivative chain rule
contains a term involving \(DS_W(\mathbf1)\), which vanishes by criticality;
the two remaining first derivatives of \(\operatorname{Exp}\) each supply
the factor \(N_{\mathrm{mesh}}^{-1}\).  This proves
\eqref{cusp-link-exponential-hessian-scaling}.
\end{proof}

The Gaussian Hilbert space in \eqref{cusp-link-gaussian-hilbert-space} is
not the compact-link Yang--Mills Hilbert space.  At this finite lattice, the
latter kinematic and gauge-invariant spaces would be
\begin{equation}
 \mathcal H^{\mathrm{link}}_{N_{\mathrm{mesh}}}
 =L^2(\mathrm{SU}(2)^{E_{N_{\mathrm{mesh}}}},dU),
 \qquad
 \mathcal H^{\mathrm{inv}}_{N_{\mathrm{mesh}}}
 =\left(\mathcal H^{\mathrm{link}}_{N_{\mathrm{mesh}}}\right)^{
    \mathrm{SU}(2)^{V_{N_{\mathrm{mesh}}}}},
 \label{cusp-link-physical-hilbert-space}
\end{equation}
where \(dU\) is product normalized Haar measure and the superscript denotes
the nonlinear vertex-gauge action.  Here ``gauge-invariant'' refers only to
the link-configuration space of this four-dimensional cell complex.  Without
a chosen time direction or transfer matrix, it is not asserted to be a
canonical Hamiltonian physical Hilbert space.  L\"uscher's construction of
a selfadjoint, strictly positive transfer matrix
\cite{Luscher_1977} and Menotti--Pelissetto's reflection-positivity theorem
for the Wilson action \cite{MenottiPelissetto_1987} concern the standard
Wilson lattice setting after the requisite reflection/time structure has
been fixed; they do not supply those omitted choices for the nonorthogonal
sector functional \eqref{cusp-link-wilson-sector-definition}.  The
Osterwalder--Schrader reconstruction boundary is likewise the one stated in
the primary Euclidean axioms papers
\cite{OsterwalderSchrader_1973,OsterwalderSchrader_1975}, not a consequence of
a finite Gaussian oscillator calculation.  The differences are
structural:
\begin{enumerate}
\item \(\mathcal H^{\mathrm{Gauss}}\) is an \(L^2\)-space on a noncompact
linear Coulomb slice with Lebesgue measure; \(\mathcal H^{\mathrm{inv}}\)
is a gauge-invariant subspace of an \(L^2\)-space on compact link variables
with Haar measure.
\item The local exponential map is neither a global coordinate chart nor a
measure-preserving bijection.  No unitary map between the two displayed
Hilbert spaces has been constructed.
\item A link configuration \(U\) determines the Dirac distribution
\(\delta_U\), but
\begin{equation}
 \delta_U\in\mathcal D'(\mathrm{SU}(2)^{E_{N_{\mathrm{mesh}}}}),
 \qquad
 \delta_U\notin L^2(\mathrm{SU}(2)^{E_{N_{\mathrm{mesh}}}},dU).
 \label{cusp-link-dirac-not-l2}
\end{equation}
Indeed, product Haar measure is nonatomic, so a point mass is singular with
respect to it and has no square-integrable density.
\item No nonlinear compact-link Hamiltonian, Gauss-law spectral problem,
time direction, reflection-positive transfer matrix, renormalized-scale
limit, or Osterwalder--Schrader reconstruction is supplied here.
\end{enumerate}
Accordingly, \eqref{cusp-link-gaussian-gap} is an exact theorem about the
finite quadratic oscillator defined by the Coulomb-restricted operator
\(g_{\mathrm{YM}}^{-2}d^*d\).  It is not an eigenvalue statement about a
nonperturbative physical Yang--Mills
Hamiltonian, and the small classical action in
\eqref{cusp-link-log-action-value} is not by itself a physical-state Rayleigh
quotient.  Similarly, \(\widetilde\Delta_{T,M_{\mathrm{cov}}}\) is a scalar
geometric Laplacian and \(\mathcal L_T^{\mathrm{cont}}\) is a classical
Yang--Mills Hessian at the product connection.  No common direct-integral
operator, physical Hamiltonian, or map from either of those operators to the
 compact-link Hilbert space is asserted.

\subsection{Wilson-loop states, entanglement, and the exact coercive bridge}
\label{cusp-link-spectral-logic}

The preceding distinction can be made by exact countermodels, rather than by
terminology.  Put
\[
 \mathscr X_N=\mathrm{SU}(2)^{E_{N_{\mathrm{mesh}}}},
 \qquad
 \mathscr G_N=\mathrm{SU}(2)^{V_{N_{\mathrm{mesh}}}},
\]
and use the vertex-gauge action compatible with
\eqref{cusp-link-inverse-transport-gauge-law}:
\begin{equation}
 (g\cdot U)_{n,j}
 =g(n)^{-1}U_{n,j}g(n+\mathbf e_j).
 \label{cusp-link-vertex-gauge-action}
\end{equation}
If \(C=(e_1^{\varepsilon_1},\ldots,e_r^{\varepsilon_r})\) is an oriented
closed edge path, with \(\varepsilon_k\in\{1,-1\}\), let
\begin{equation}
 W_C(U)=\frac12\operatorname{Re}\operatorname{tr}
 \left(U_{e_1}^{\varepsilon_1}\cdots U_{e_r}^{\varepsilon_r}\right).
 \label{cusp-link-wilson-loop-observable}
\end{equation}
Here an edge with exponent \(-1\) is traversed opposite to its chosen
orientation.  Under \eqref{cusp-link-vertex-gauge-action}, the ordered product
around \(C\) is conjugated at its base vertex, so \(W_C\) is gauge invariant.
A \emph{Euclidean link state} in this paragraph means a
\(\mathscr G_N\)-invariant Borel probability measure on \(\mathscr X_N\).
This is deliberately not called a Hamiltonian state: no time direction or
transfer matrix has been chosen.

\begin{proposition}[Link kinematics and one magnetic field select no area law]
\label{cusp-link-no-area-law-from-kinematics}
There are \(\mathscr G_N\)-invariant Borel probability measures
\(\mu_{\mathrm{pg}}\) and \(\mu_{\mathrm H}\) on the same configuration
space \(\mathscr X_N\) such that
\begin{align}
 \int_{\mathscr X_N}W_C\,d\mu_{\mathrm{pg}}&=1
 &&\text{for every closed edge path }C,
 \label{cusp-link-pure-gauge-wilson}\\
 \int_{\mathscr X_N}W_C\,d\mu_{\mathrm H}&=0
 &&\text{if some unoriented edge occurs exactly once in }C.
 \label{cusp-link-haar-wilson}
\end{align}
Consequently, for each \(0\le t\le1\), the invariant state
\(\mu_t=t\mu_{\mathrm{pg}}+(1-t)\mu_{\mathrm H}\) has expectation \(t\)
for every path covered by the second line.  Thus the configuration space,
gauge action, and existence of the explicit flux in
Theorem~\ref{cusp-link-flux-theorem} do not determine Wilson-loop decay or a
string tension.
\end{proposition}

\begin{proof}
Let \(dg_V\) be normalized product Haar measure on \(\mathscr G_N\), and let
\(\mu_{\mathrm{pg}}\) be its push-forward under
\[
 \mathscr G_N\longrightarrow\mathscr X_N,
 \qquad
 g\longmapsto g\cdot\mathbf1.
\]
Bi-invariance of \(dg_V\) makes this push-forward gauge invariant: applying a
fixed gauge transformation replaces \(g\) by its pointwise right translate.
Along a
positively oriented edge \(v\to w\), the corresponding link is
\(g(v)^{-1}g(w)\); the inverse occurs for the opposite orientation.  The
factors telescope around every closed path, giving the identity matrix and
proving \eqref{cusp-link-pure-gauge-wilson}.

Let \(\mu_{\mathrm H}\) be normalized product Haar measure on
\(\mathscr X_N\).  Every fixed gauge transformation acts on each link by a
left and a right translation, so bi-invariance of Haar measure proves gauge
invariance.  Suppose an unoriented edge \(e\) occurs exactly once in \(C\).
After conditioning on all other links, the matrix inside the trace in
\eqref{cusp-link-wilson-loop-observable} is of the form \(AU_eB\) or
\(AU_e^{-1}B\), with \(A,B\) fixed.  In the defining representation,
\[
 \int_{\mathrm{SU}(2)}U\,dU=0,
 \qquad
 \int_{\mathrm{SU}(2)}U^{-1}\,dU=0:
\]
left translation by the central element \(-I\) changes either integral to its
negative.  Fubini's theorem proves \eqref{cusp-link-haar-wilson}.  Linearity of
integration gives the assertion for \(\mu_t\).  All these measures exist
independently of the singled-out point of \(\mathscr X_N\) defined in
\eqref{cusp-link-flux-u1}--\eqref{cusp-link-flux-u34}.
\end{proof}

The word \emph{entangled} also requires a specified tensor product.  For
complex Hilbert spaces \(\mathfrak h_A,\mathfrak h_B\) of dimension at least
two, a nonzero vector of \(\mathfrak h_A\otimes\mathfrak h_B\) is a product
vector if it equals \(a\otimes b\); otherwise it is entangled.

\begin{proposition}[Entanglement is neither a sector nor an energy floor]
\label{cusp-link-entanglement-no-gap}
The entangled vectors, together with zero, do not form a linear subspace.
They are norm dense and have dense linear span, so no orthogonal projection
has precisely the entangled vectors as its range.  Moreover:
\begin{enumerate}
\item on \(\mathbb C^2\otimes\mathbb C^2\), the bounded nonnegative
self-adjoint operator
\[
 H_B=I-|\Omega_B\rangle\langle\Omega_B|,
 \qquad
 \Omega_B=2^{-1/2}(|00\rangle+|11\rangle),
\]
with domain the whole tensor product, has the entangled vector \(\Omega_B\)
as its unique zero-energy ground state and has gap one;
\item there is a bounded nonnegative self-adjoint operator, again with full
domain, which has a unique entangled zero-energy ground state and entangled
positive-energy eigenvectors whose eigenvalues converge to zero.
\end{enumerate}
Hence a statement that some degrees of freedom must be entangled supplies
neither a positive energy nor a positive spectral gap.
\end{proposition}

\begin{proof}
Choose orthonormal \(e_0,e_1\in\mathfrak h_A\) and
\(f_0,f_1\in\mathfrak h_B\).  The Bell vectors
\[
 \Psi_\pm=2^{-1/2}(e_0\otimes f_0\pm e_1\otimes f_1)
\]
have coefficient-matrix rank two and are entangled, whereas their sum is the
product vector \(\sqrt2\,e_0\otimes f_0\).  Thus the set is not linear.  Any
product vector \(a\otimes b\) is the norm limit, as \(\varepsilon\to0\), of
\(a\otimes b+\varepsilon a'\otimes b'\), where \(a'\perp a\) and
\(b'\perp b\); for \(\varepsilon\ne0\) the latter has coefficient rank two.
This proves density.  Also
\[
 a\otimes b=\frac12\bigl[(a\otimes b+a'\otimes b')
                         +(a\otimes b-a'\otimes b')\bigr],
\]
so product vectors, which span a dense subspace, lie in the span of entangled
vectors.  A projection range is a closed linear subspace.

For the first operator assertion, the rank-one term is the orthogonal
projection onto the Bell line.  Hence \(H_B\) has spectrum \(\{0,1\}\) and
kernel \(\mathbb C\Omega_B\).

For the second, take
\(\mathfrak h_A=\mathfrak h_B=\ell^2(\mathbb N_0)\), with orthonormal bases
\((e_j)\) and \((f_j)\), and set
\[
 \Omega=2^{-1/2}(e_0\otimes f_0+e_1\otimes f_1),
 \qquad
 \Phi_n=2^{-1/2}(e_{2n}\otimes f_{2n+1}
                         +e_{2n+1}\otimes f_{2n})\quad(n\ge1).
\]
These vectors are orthonormal.  Define \(H\), on the whole tensor product, by
\begin{equation}
 H\Omega=0,
 \qquad H\Phi_n=n^{-1}\Phi_n,
 \qquad H=I\quad\text{on }
 \{\Omega,\Phi_1,\Phi_2,\ldots\}^{\perp}.
 \label{cusp-link-entangled-gapless-operator}
\end{equation}
This real diagonal prescription defines a bounded nonnegative self-adjoint
operator.  Only \(\Omega\) has eigenvalue zero.  Each \(\Phi_n\) is entangled
and has eigenvalue \(1/n\), which tends to zero.  More strongly, every
nonzero vector in
\(\overline{\operatorname{span}}\{\Phi_n:n\ge1\}\) is entangled: its
coefficient operator has a nonzero rank-two block
\(2^{-1/2}c_n\left(\begin{smallmatrix}0&1\\1&0\end{smallmatrix}\right)\)
for some \(n\).
\end{proof}

The same logical obstruction occurs even without a tensor factorization.
On \(\ell^2(\mathbb N_0)\), let the Gaussian number operator be
\(N|n\rangle=n|n\rangle\) on
\(\operatorname{Dom}(N)=\{(c_n):\sum n^2|c_n|^2<\infty\}\).  It has gap one.
The bounded nonnegative self-adjoint operator \(F(N)\), defined on the whole
space by
\begin{equation}
 F(N)|0\rangle=0,
 \qquad F(N)|1\rangle=|1\rangle,
 \qquad F(N)|n\rangle=n^{-1}|n\rangle\quad(n\ge2),
 \label{cusp-link-same-one-particle-gapless-completion}
\end{equation}
agrees with \(N\) on the vacuum and the entire one-particle sector, yet has
unique vacuum and zero excitation gap.  This follows directly from the
displayed eigenvalues.  Thus even exact vacuum and one-particle Gaussian data
do not determine a nonlinear completion.

We finally record the positive statement that identifies the missing map.
Let \(A_{\mathrm{phys}}\) be a lower-semibounded self-adjoint operator on a
Hilbert space \(\mathcal H_{\mathrm{phys}}\), with ground energy \(E_0\),
nonzero ground projection
\[
 P_0=\mathbf1_{\{E_0\}}(A_{\mathrm{phys}}),
 \qquad
 \mathcal H_{\mathrm{exc}}=(I-P_0)\mathcal H_{\mathrm{phys}},
\]
closed form \(\mathfrak a_{\mathrm{phys}}\), and form domain
\(\mathcal Q(A_{\mathrm{phys}})\).  Let \(B\ge0\) be an auxiliary
self-adjoint operator with closed form \(\mathfrak b\), form domain
\(\mathcal Q(B)\), and \(\mathfrak b[v]\ge\gamma\|v\|^2\) for some
\(\gamma>0\).

\begin{theorem}[Full-domain coercive gap transfer]
\label{cusp-link-full-domain-gap-transfer}
Suppose a linear map
\[
 R:\mathcal Q(A_{\mathrm{phys}})\cap\mathcal H_{\mathrm{exc}}
   \longrightarrow\mathcal Q(B)
\]
is defined on the entire physical excitation form domain and, for constants
\(c,C>0\), satisfies
\begin{equation}
 \|Ru\|\ge c\|u\|,
 \qquad
 \mathfrak b[Ru]\le
 C\bigl(\mathfrak a_{\mathrm{phys}}[u]-E_0\|u\|^2\bigr).
 \label{cusp-link-full-domain-comparison}
\end{equation}
Then
\begin{equation}
 \inf\operatorname{Spec}
 \left(A_{\mathrm{phys}}|_{\mathcal H_{\mathrm{exc}}}\right)-E_0
 \ge\frac{\gamma c^2}{C}>0.
 \label{cusp-link-transferred-gap}
\end{equation}
By contrast, an isometric intertwining of \(B\) with only one reducing
subspace of \(A_{\mathrm{phys}}\) does not imply a full gap.
\end{theorem}

\begin{proof}
For every \(u\) in the stated physical form domain,
\[
 C\bigl(\mathfrak a_{\mathrm{phys}}[u]-E_0\|u\|^2\bigr)
 \ge\mathfrak b[Ru]
 \ge\gamma\|Ru\|^2
 \ge\gamma c^2\|u\|^2.
\]
The ground spectral projection reduces \(A_{\mathrm{phys}}\).  The spectral
theorem, applied to its excitation part, identifies the infimum of the
Rayleigh quotient of the shifted closed form with the left-hand side of
\eqref{cusp-link-transferred-gap}, proving the bound.

For the final assertion, take \(B=I\) on \(\mathbb C\), and on
\[
 \mathcal H=\mathbb C\Omega\oplus\mathbb C\oplus L^2((0,1),dx)
\]
take the bounded nonnegative self-adjoint operator
\[
 A=0\oplus I\oplus M_x,
\]
where \((M_xf)(x)=xf(x)\).  The isometry
\(Jz=0\oplus z\oplus0\) obeys \(AJ=JB\), and its range is reducing and
gapped.  Nevertheless \(A\) has unique vacuum \(\Omega\), while normalized
functions supported in \((0,1/n)\) have Rayleigh quotient at most \(1/n\).
Thus the full excitation gap is zero.
\end{proof}

In the present paper, the auxiliary Gaussian oscillator is explicit, but no
map \(R\) satisfying \eqref{cusp-link-full-domain-comparison} from a physical
compact-link excitation form domain has been constructed.  Moreover, its own
gap tends to zero along \eqref{cusp-link-diagonal-sequence}; hence it cannot
provide a uniform positive \(\gamma\) on that sequence.  The Hilbert space
\eqref{cusp-link-physical-hilbert-space} also contains link variables only:
no fundamental-matter Hilbert factor, pair-creation operator, hopping term,
or matter mass term has been defined.  String breaking by dynamical matter is
therefore not an observable or spectral process typed by this model.

\section{Conclusion}

The constructions above give an exact compatibility diagram at the level
that has actually been proved:
\begin{align*}
 \{\text{smooth connections on }\widetilde F_{T,M_{\mathrm{cov}}}\}
 &\xrightarrow{\ \operatorname{Hol}\ }
 \mathrm{SU}(2)^{E_{N_{\mathrm{mesh}}}},\\
 \mathscr C_{T,M_{\mathrm{cov}},N_{\mathrm{mesh}}}
 &\xrightarrow[\text{locally}]{\ \operatorname{Exp}_{N_{\mathrm{mesh}}}\ }
 \mathrm{SU}(2)^{E_{N_{\mathrm{mesh}}}}.
\end{align*}
The first arrow is globally defined but is not claimed to be surjective.
The second is a local coordinate map at zero, not a global embedding and not
a measure-preserving identification.  On the compact-link side, the
 periodic flux and both displayed action values are exact.  On the linear
 side, the scalar cover spectrum, deck-invariant sector, continuum native
 Hessian spectrum, finite Fourier eigenvalues, and oscillator gap are exact.
 On the distinct smooth native side, the unit-flux line bundle, the reducible
 \(\mathrm{SU}(2)\) Yang--Mills critical connection, and its action
 \(2\pi^2/(g_{\mathrm{YM}}^2D_T)\) are exact; the explicit charge-two theta
 section produces the reduced Jacobi eigenvalue \(-4\pi/D_T\).  Thus that
 smooth magnetic representative is a saddle at every finite \(T\), and its
 numerical action agreement with the native link field does not identify the
 two configurations.
 The native scalar direct integral is also an exact self-adjoint operator with
trivial kernel and zero in its essential spectrum.
The native torus tends pointedly to \(\mathbb T^2\times\mathbb R^2\), whereas
every sequence with \(T_j\to\infty\) and
\(M_{\mathrm{cov},j}\to\infty\) tends pointedly to \(\mathbb R^4\).  Along the
explicit diagonal sequence, the cover systole diverges, physical mesh edges
shrink, and the Gaussian gap and logarithmic action vanish.

The same link kinematics admits gauge-invariant Euclidean states whose simple
Wilson-loop expectations are respectively one, zero, or any prescribed
constant between them.  The magnetic configuration therefore selects no
confinement law.  Likewise, entanglement does not define a spectral subspace
and imposes no energy floor, while exact agreement with one gapped Gaussian
range leaves open a gapless orthogonal channel.  Theorem~\ref{cusp-link-full-domain-gap-transfer}
isolates the sufficient missing datum: a norm-nondegenerate energy comparison
defined on the entire physical excitation form domain.

Nothing in this diagram identifies the finite Gaussian Hilbert space with a
physical gauge-invariant compact-link Hilbert space.  In particular, the
vanishing upper bound for the Gaussian gap is a theorem about the explicitly
defined quadratic model, not a conclusion about the spectrum of a
nonperturbative Yang--Mills Hamiltonian.  This non-identification is part of
the mathematical result, not a deferred step.

\paragraph{Provenance.}
The compact-link core was extracted from the canonical consolidated
workbench source \texttt{cusp\_compact\_link\_gaussian.tex}.  The fibrewise
cover spectrum and pointed-limit results, and the native continuum and finite
Coulomb comparisons, were synchronized from the proved workbench sources
 \texttt{cusp\_spectral\_geometry.tex} and
 \texttt{cusp\_gauge\_operators.tex} and restated here with their required
 definitions and proofs.  The smooth integral-flux construction, closed
 background-gauge domains, and charged negative mode were synchronized from
 \texttt{cusp\_su2\_magnetic.tex}; they are restated here rather than invoked
 as an external lemma.  The operator-theoretic countermodels and coercive
comparison in \S\ref{cusp-link-spectral-logic} were cross-checked against the
internal exact audit module
\path{agent_logs/nested_sector_gap_theorem_20260903.tex} and are
reproved here.  The workbench sources themselves are not modified.

\begingroup
\small
\bibliographystyle{plain}
\bibliography{references}
\endgroup

\end{document}
